\documentclass[11pt,a4paper]{article}
\usepackage[utf8]{inputenc}
\usepackage[margin=1in]{geometry}
\usepackage{amsmath}
\usepackage{amssymb}
\usepackage{enumitem}
\usepackage{booktabs}
\usepackage{titlesec}
\usepackage{microtype}
\usepackage{xcolor}

\titleformat{\section}{\large\bfseries\sffamily}{}{0em}{}[\titlerule]
\titleformat{\subsection}{\normalsize\bfseries\sffamily}{}{0em}{}

\setlength{\parindent}{0pt}
\setlength{\parskip}{7pt plus 2pt minus 1pt}

\title{\textbf{CAMS Comprehensive Hard Exam --- Full Topic Coverage}}
\author{\textbf{Advanced AML Training Framework}}
\date{\today}

\begin{document}
	
	\maketitle
	
	\section*{Instructions}
	This exam contains \textbf{70 questions} across all CAMS exam domains. Each question has \textbf{five answer choices (A--E)}.
	\begin{itemize}
		\item Questions marked \textbf{[SELECT ONE]} have exactly one correct answer.
		\item Questions marked \textbf{[SELECT TWO]} have exactly two correct answers. You must select both to receive credit.
	\end{itemize}
	Questions are scenario-based. Read each scenario carefully before answering. Partial credit is not awarded on SELECT TWO questions.
	
	\newpage
	%==============================================================
	\section{Part 1: Money Laundering Methods, Stages, and Typologies}
	%==============================================================
	
	%--- Q1 ---
	\subsection{1. Three-Stage Analysis in a Layered Cash Scheme \textnormal{[SELECT ONE]}}
	
	An organized crime group controlling a network of nail salons in three U.S. states generates approximately \$4 million per year in illicit narcotics proceeds. The group's accountant deposits cash at 14 different branch locations of four separate banks, deliberately keeping each deposit between \$7,500 and \$9,800. The deposited funds are then consolidated into a holding LLC, which issues fictitious management-fee invoices to all 14 salons. The LLC uses the aggregated funds to purchase three residential rental properties. Rental income from the properties is reported to the IRS as legitimate business income. A transaction monitoring analyst at one of the banks flags the nail salon deposits.
	
	Which stage of money laundering most accurately describes the act of issuing fictitious management-fee invoices to consolidate funds in the LLC?
	
	\begin{enumerate}[label=\Alph*.]
		\item Placement, because the invoices introduce criminal proceeds into the financial system for the first time.
		\item Layering, because the invoices create a fictitious paper trail that obscures the criminal origin of the deposited cash by generating an apparently legitimate business explanation.
		\item Integration, because the funds are now documented as business expenses deductible against salon revenue.
		\item Structuring, because the invoices are designed to keep each transaction below reporting thresholds.
		\item Predicate offense, because invoice fraud is itself a standalone financial crime.
	\end{enumerate}
	\textbf{Correct Answer: B}
	
	%--- Q2 ---
	\subsection{2. Trade-Based Money Laundering: Over-Invoicing vs. Under-Invoicing \textnormal{[SELECT ONE]}}
	
	A compliance officer at a trade finance bank reviews a letter of credit supporting the import of generic pharmaceutical raw materials from an exporter in Country A (a jurisdiction with high corruption indices) to an importer in Country B (a FATF-compliant country). The unit price on the invoice is \$1,200 per kilogram. An independent commodity pricing database confirms that the international market price for the same raw material is \$180 per kilogram. The importer in Country B is a licensed pharmaceutical manufacturer with a clean compliance history. The exporter in Country A was incorporated 11 days ago and has no verifiable manufacturing infrastructure. Payment is routed through the exporter's account at a shell bank in a secrecy haven.
	
	What is the primary TBML technique being used, and what is the underlying criminal purpose?
	
	\begin{enumerate}[label=\Alph*.]
		\item Under-invoicing by the exporter, enabling the importer to sell goods at market price and retain an untaxed cash surplus in Country B.
		\item Over-invoicing by the exporter, enabling the transfer of excess value from Country B to Country A disguised as a legitimate commercial payment, thereby moving illicit funds from a FATF-compliant jurisdiction to a high-risk secrecy haven.
		\item Phantom shipping, because no goods are physically transferred and the entire transaction is a fabricated paper exercise.
		\item Multiple invoicing, where the same goods have been invoiced more than once to different buyers to generate duplicate payments.
		\item Mirror trading, where the importer uses the same invoice to purchase matching derivatives positions on a regulated commodities exchange.
	\end{enumerate}
	\textbf{Correct Answer: B}
	
	%--- Q3 ---
	\subsection{3. The Black Market Peso Exchange: Identifying the Loop \textnormal{[SELECT ONE]}}
	
	A FinCEN analysis report flags a cluster of suspicious activity involving a U.S. textile importer in Los Angeles. Over 18 months, the importer's bank accounts receive structured cash deposits totaling \$8.7 million from 43 different individuals, none of whom have any apparent business relationship with the importer. Concurrently, a Colombian textile wholesaler in Medell\'in receives shipments of finished garments at prices 35\% below the standard domestic wholesale price, paid in Colombian pesos by a local broker who has no documented relationship with the U.S. importer. No funds are wired directly between Colombia and the United States.
	
	Which of the following most precisely explains the criminal mechanism at work?
	
	\begin{enumerate}[label=\Alph*.]
		\item The Colombian broker is using drug cartel dollars deposited in U.S. bank accounts to pay the U.S. textile importer's commercial invoices on behalf of the Colombian wholesaler. The Colombian wholesaler repays the broker in pesos at a favorable exchange rate, converting narco-dollars into domestic pesos without any cross-border currency movement.
		\item The U.S. textile importer is structuring cash to avoid CTR reporting and depositing the proceeds into an offshore account via wire transfer, which is a standard layering technique.
		\item The Colombian wholesaler is over-invoicing the shipment of garments to justify large wire transfers from Colombia to the United States under the guise of trade finance.
		\item The Colombian broker is a licensed money remitter operating a legitimate currency exchange business in compliance with FATF Recommendation 14.
		\item The scheme is best characterized as a hawala arrangement, where the U.S. deposits and Colombian peso payments are informal remittances with no commercial trade component.
	\end{enumerate}
	\textbf{Correct Answer: A}
	
	%--- Q4 ---
	\subsection{4. Free Trade Zone Vulnerabilities and Phantom Shipments \textnormal{[SELECT TWO]}}
	
	A logistics and trade compliance officer at a major European port is investigating a shipping company that operates exclusively within the Hamburg Free Trade Zone. The company has processed 340 container shipments over 24 months, all of which are declared as ``industrial parts and components'' with invoice values uniformly listed at \$30,000 per container regardless of the listed cargo type. Physical inspection of four containers selected at random reveals that one contained the declared industrial parts, one contained household furniture, one was completely empty, and one contained luxury consumer electronics with a retail value exceeding \$400,000. Bills of lading for all containers have been altered to reflect the same generic cargo description. The shipping company's sole director is a legal nominee.
	
	Which two of the following findings most directly constitute money laundering red flags that must be escalated?
	
	\begin{enumerate}[label=\Alph*.]
		\item The uniform invoice value of \$30,000 per container regardless of actual cargo type is a hallmark of trade-based money laundering through deliberate price manipulation and mislabeling.
		\item Shipping companies operating in free trade zones are exempt from AML scrutiny because FTZs are supervised exclusively by customs authorities.
		\item The presence of an empty container among the sampled shipments suggests phantom shipping, where false documentation is generated without any physical movement of goods of commercial value.
		\item The nominee director structure is a standard corporate governance practice acceptable under European company law with no AML implications.
		\item The Hamburg Free Trade Zone applies the same AML monitoring standards as any regular commercial port.
	\end{enumerate}
	\textbf{Correct Answers: A and C}
	
	%--- Q5 ---
	\subsection{5. Prepaid Cards: Cross-Border Currency Smuggling Alternative \textnormal{[SELECT ONE]}}
	
	Customs agents at a major international airport stop a traveler carrying 62 open-loop prepaid debit cards, each loaded with \$9,500, purchased at various pharmacy and supermarket retail locations across five U.S. cities over the preceding 30 days. The traveler is unable to explain the source of the funds or provide a business purpose for carrying 62 cards. The cards were purchased with cash and are not linked to any bank account. The traveler is flying to a country under FinCEN heightened scrutiny for money laundering. The total value across all cards is \$589,000.
	
	What is the primary AML concern presented by this scenario, and what regulatory principle does it implicate?
	
	\begin{enumerate}[label=\Alph*.]
		\item The traveler may have violated structuring laws by purchasing each card at exactly \$9,500 to remain below the \$10,000 CTR threshold, and may have evaded currency reporting requirements by using prepaid cards, which are classified as monetary instruments under BSA regulations, to transport the equivalent of \$589,000 in currency.
		\item The cards are retail products and are not subject to any currency controls, regardless of value or number, because they are not physical currency.
		\item The only violation is a state consumer protection law violation for purchasing prepaid cards without a business license.
		\item The situation constitutes proliferation financing because the destination country is under FATF scrutiny.
		\item The customs agency should simply confiscate the cards and release the traveler, as there is no AML reporting obligation.
	\end{enumerate}
	\textbf{Correct Answer: A}
	
	%--- Q6 ---
	\subsection{6. Correspondent Banking: Nested Transactions and Liability \textnormal{[SELECT ONE]}}
	
	PrimeClear Bank (EU) maintains a USD correspondent account for Atlantic Regional Bank (ARB), a commercial bank in a Caribbean jurisdiction rated ``Partially Compliant'' on the most recent FATF mutual evaluation. PrimeClear's annual due diligence review of ARB's correspondent account reveals that ARB's account is being used to route payments for CashPoint FX, an unlicensed currency exchange operating in the same jurisdiction. CashPoint FX is not disclosed in any ARB documentation provided to PrimeClear. The transactions flowing through the nested arrangement include large round-dollar wires totaling \$23 million over six months, with no identifiable commercial purpose stated. PrimeClear's compliance director argues that because ARB holds a valid domestic banking license, PrimeClear has no residual obligation for transactions originated by ARB's undisclosed downstream clients.
	
	Which of the following statements most accurately assesses PrimeClear's legal and regulatory position?
	
	\begin{enumerate}[label=\Alph*.]
		\item PrimeClear's compliance director is correct. The existence of a valid domestic banking license held by ARB transfers full AML responsibility to ARB under the principle of respondent bank primary accountability.
		\item PrimeClear faces direct regulatory exposure. Correspondent banks are responsible for understanding who is actually accessing their payment infrastructure. The undisclosed nesting of an unlicensed currency exchange into the correspondent account means PrimeClear has lost visibility into the true originators and beneficiaries, a condition that violates core KYC and correspondent banking due diligence standards. PrimeClear must issue a Request for Information to ARB, file an STR if warranted, and reassess whether the relationship should continue.
		\item PrimeClear's only obligation is to lower ARB's risk tier on its internal risk matrix and increase transaction monitoring threshold sensitivity for ARB's account.
		\item Nested transactions are a routine feature of correspondent banking, and PrimeClear's exposure is limited to filing a CTR for any single wire exceeding \$10,000.
		\item PrimeClear is protected from regulatory action because it did not have knowledge of the nesting arrangement and therefore cannot be held responsible.
	\end{enumerate}
	\textbf{Correct Answer: B}
	
	%--- Q7 ---
	\subsection{7. Digital Assets: Layering Through a DeFi Protocol \textnormal{[SELECT TWO]}}
	
	A blockchain analytics firm retained by a regulated crypto exchange identifies a suspicious activity pattern involving a customer, Wallet Alpha. Over 48 hours, Wallet Alpha receives 200 Ethereum (ETH) from a wallet flagged by OFAC as associated with ransomware proceeds. Within minutes, the 200 ETH is deposited into a decentralized lending protocol where it is used as collateral for a \$320,000 stablecoin (USDC) loan. The USDC is immediately bridged to a second blockchain using a cross-chain bridge, converted to a privacy token (XMR), routed through a mixing protocol, and then bridged back to the original blockchain as WETH (wrapped Ether) two weeks later. The WETH is deposited into the registered exchange with a customer claiming the funds are ``yield farming proceeds.''
	
	Which two of the following represent the most critical findings requiring immediate action by the exchange?
	
	\begin{enumerate}[label=\Alph*.]
		\item The funds deposited at the exchange are directly traceable through on-chain analysis to an OFAC-flagged wallet, triggering mandatory sanctions screening and asset blocking obligations.
		\item The use of a cross-chain bridge, privacy token conversion, and mixing protocol constitutes a multi-stage layering sequence specifically designed to sever the transactional link to the ransomware-associated source wallet, requiring the exchange to file a SAR and block the withdrawal.
		\item The exchange has no obligation to investigate because yield farming is a legitimate decentralized finance activity protected under financial innovation safe harbors.
		\item The funds are now clean because they have passed through multiple blockchain protocols and the original flagged address is no longer directly visible.
		\item The exchange's only obligation is to update the customer's risk profile to ``medium risk'' and apply standard enhanced due diligence at the next annual review.
	\end{enumerate}
	\textbf{Correct Answers: A and B}
	
	%==============================================================
	\section{Part 2: FATF Standards, Mutual Evaluations, and International Frameworks}
	%==============================================================
	
	%--- Q8 ---
	\subsection{8. FATF Mutual Evaluation: Technical Compliance vs.\ Effectiveness \textnormal{[SELECT ONE]}}
	
	Country Meridian has enacted a comprehensive AML/CFT legal framework that maps directly to all 40 FATF Recommendations. The country has a functioning Financial Intelligence Unit that receives and processes over 30,000 STRs annually. All major financial institutions have implemented automated transaction monitoring, sanctions screening, and PEP identification systems. However, a FATF Mutual Evaluation Assessment Team notes the following: the FIU's financial intelligence has never been disseminated to law enforcement for use in a criminal investigation; the country has secured zero asset forfeitures under its AML law despite having enacted the relevant legislation four years ago; and domestic courts have not prosecuted a single stand-alone money laundering case in seven years, citing a lack of training among prosecutors on financial crime methodology.
	
	How will Country Meridian's Mutual Evaluation most likely be rated across both assessment pillars?
	
	\begin{enumerate}[label=\Alph*.]
		\item Largely Compliant on Technical Compliance (laws are in place); Low Effectiveness on multiple Immediate Outcomes, particularly IO.6 (financial intelligence used by competent authorities), IO.7 (ML investigated and prosecuted), and IO.8 (proceeds of crime confiscated).
		\item Non-Compliant on Technical Compliance because the failure to prosecute demonstrates that the legislation does not actually meet FATF standards; Low Effectiveness on all Immediate Outcomes.
		\item Compliant on Technical Compliance and Largely Effective on Effectiveness because the volume of STRs received demonstrates system activity.
		\item Partially Compliant on Technical Compliance because the legislation has not been effectively tested in courts; Moderate Effectiveness because STR volumes are high.
		\item Compliant on both pillars, because FATF Mutual Evaluations do not consider prosecution outcomes.
	\end{enumerate}
	\textbf{Correct Answer: A}
	
	%--- Q9 ---
	\subsection{9. FATF Grey List vs.\ Black List: Institutional Obligations \textnormal{[SELECT TWO]}}
	
	A risk committee at a multinational bank is reviewing its exposure to two jurisdictions. Jurisdiction Alpha has just been placed on FATF's ``Increased Monitoring'' list (Grey List) because of significant deficiencies in its AML/CFT framework, though it has committed to an action plan. Jurisdiction Beta remains on FATF's ``High-Risk Jurisdictions subject to a Call for Action'' (Black List) for strategic deficiencies, and FATF has called on its members to apply countermeasures. The bank currently maintains correspondent accounts with one bank in each jurisdiction and provides trade finance services to exporters in both countries.
	
	Which two statements correctly describe the bank's obligations with respect to these two jurisdictions?
	
	\begin{enumerate}[label=\Alph*.]
		\item For Jurisdiction Alpha (Grey List), the bank must apply Enhanced Due Diligence to all transactions and correspondent accounts but is not automatically required to terminate the relationships.
		\item For Jurisdiction Beta (Black List), FATF's call for countermeasures means the bank should evaluate and implement appropriate countermeasures, which may include severely restricting or terminating transactions and correspondent relationships with entities from that jurisdiction.
		\item Both jurisdictions require identical treatment: all transactions must be frozen and all accounts closed immediately.
		\item The bank has no obligation to take any specific action for Grey List jurisdictions because the listing is advisory only.
		\item Countermeasures against Black List jurisdictions are optional for private financial institutions; only central banks are obligated to act.
	\end{enumerate}
	\textbf{Correct Answers: A and B}
	
	%--- Q10 ---
	\subsection{10. FATF Recommendation 16: Wire Transfer Information \textnormal{[SELECT ONE]}}
	
	An international payments processor, SwiftRoute Ltd., is reviewing its compliance with FATF Recommendation 16 (the Wire Transfer Rule). A spot-check audit of 200 recent cross-border wire transfers reveals the following deficiency patterns: (1) 47 transfers above USD 1,000 contained only the originator's name and account number but omitted the originator's address, national identification number, or date of birth; (2) 18 transfers contained the originator's full required information but omitted the beneficiary's name and account number; and (3) 22 transfers were below the USD 1,000 threshold and contained only the originator's name and account number without a full address.
	
	Which category of deficiency represents the most critical regulatory violation requiring immediate remediation?
	
	\begin{enumerate}[label=\Alph*.]
		\item Category (3), because even sub-threshold transfers must carry full originator information under FATF Recommendation 16.
		\item Category (1), because transfers above USD 1,000 must contain the originator's name, account number, and at least one of: address, national ID number, or date of birth. Omitting all three qualifying identifiers on 47 transfers is a direct and quantifiable violation of the Recommendation 16 full originator information requirement.
		\item Category (2), because the beneficiary's information is required for all cross-border transfers above USD 1,000 under the Travel Rule, and 18 transfers with missing beneficiary information exposes the receiving institution to sanctions screening failures.
		\item Categories (1) and (2) are equally critical; Category (3) is not a violation.
		\item None of the categories represent a violation because SwiftRoute processes transactions rather than originating them, and intermediary institutions have reduced obligations under Recommendation 16.
	\end{enumerate}
	\textbf{Correct Answer: B}
	
	%--- Q11 ---
	\subsection{11. EU 6AMLD: Predicate Offenses and Corporate Criminal Liability \textnormal{[SELECT ONE]}}
	
	A European bank's legal team is reviewing whether a recent enforcement action taken by a domestic regulator against a corporate subsidiary triggers liability under the EU Sixth Anti-Money Laundering Directive (6AMLD). The subsidiary, a trading company, was found to have knowingly processed 47 wire transfers totaling \$12 million in proceeds derived from tax fraud, which is designated as a predicate offense under 6AMLD. The country's domestic 6AMLD implementing legislation explicitly extends criminal liability for money laundering to legal persons (corporations), not just natural persons.
	
	Which of the following most accurately describes the subsidiary's exposure under 6AMLD?
	
	\begin{enumerate}[label=\Alph*.]
		\item The subsidiary faces no criminal exposure because 6AMLD only harmonizes predicate offenses for natural persons, not corporate entities.
		\item The subsidiary faces criminal liability for money laundering as a legal person under 6AMLD. The Directive harmonizes 22 predicate offenses including tax crimes, and explicitly extends criminal liability to legal persons, meaning the corporation itself, not only its individual officers, can be prosecuted and subject to criminal sanctions such as fines, debarment from public contracts, and dissolution.
		\item The subsidiary is only liable if the individual directors are also convicted; corporate liability cannot exist independently of individual liability under 6AMLD.
		\item Tax fraud was not included in the 6AMLD harmonized list of predicate offenses and therefore cannot give rise to a money laundering charge.
		\item 6AMLD applies only to financial institutions regulated by national banking supervisors, not to trading companies.
	\end{enumerate}
	\textbf{Correct Answer: B}
	
	%--- Q12 ---
	\subsection{12. USA PATRIOT Act Sections 311 and 319(b): Shell Banks and Interbank Accounts \textnormal{[SELECT TWO]}}
	
	GlobalClear Bank, a U.S. federally chartered bank, has maintained a correspondent account for the past three years on behalf of Pacific Offshore Bank, a financial institution licensed in a Pacific Island jurisdiction. GlobalClear's annual due diligence review reveals that Pacific Offshore Bank has no physical office in any jurisdiction, employs no staff, conducts all operations through a shared services agreement with an unrelated corporate services firm, and maintains no regulatory reporting obligations in its country of license because of an exemption for ``shell holding company banks'' in the local banking law. Pacific Offshore has processed \$340 million in wire transfers through the GlobalClear account in the past 12 months.
	
	Which two of the following actions are mandatory for GlobalClear under the USA PATRIOT Act?
	
	\begin{enumerate}[label=\Alph*.]
		\item GlobalClear must terminate the correspondent account immediately, because Section 313 of the USA PATRIOT Act absolutely prohibits U.S. financial institutions from maintaining correspondent accounts for foreign shell banks.
		\item If Pacific Offshore Bank cannot demonstrate that it is affiliated with a regulated, physically present depository institution that is subject to comprehensive AML supervision, GlobalClear must close the account without exception.
		\item GlobalClear may continue the relationship as long as it files a SAR every 90 days documenting the shell bank risk.
		\item Section 319(b) of the USA PATRIOT Act provides the legal mechanism for GlobalClear to demand the production of records from Pacific Offshore's U.S. correspondent account to satisfy law enforcement subpoenas, and allows for interbank account seizure if the bank refuses to comply.
		\item GlobalClear should reclassify Pacific Offshore as a ``low-risk shell entity'' and reduce its transaction monitoring to semi-annual reviews.
	\end{enumerate}
	\textbf{Correct Answers: A and D}
	
	%--- Q13 ---
	\subsection{13. Wolfsberg Group: Correspondent Banking Due Diligence Questionnaire \textnormal{[SELECT ONE]}}
	
	A global correspondent bank, NorthHarbor Bank, is onboarding a new respondent bank, SunCoast Commercial Bank, headquartered in a jurisdiction with a moderate FATF rating. As part of its due diligence, NorthHarbor receives SunCoast's Wolfsberg Correspondent Banking Due Diligence Questionnaire (CBDDQ). The CBDDQ reveals the following: SunCoast services 14 Money Services Business customers; three of these MSBs offer unlicensed international remittance services; SunCoast provides payable-through account access to 22 sub-agents; and SunCoast's AML program was last independently audited four years ago with no audit scheduled for the foreseeable future. SunCoast's local regulator has issued no formal enforcement actions.
	
	According to Wolfsberg principles, which factor in SunCoast's CBDDQ profile represents the highest residual AML risk to NorthHarbor Bank?
	
	\begin{enumerate}[label=\Alph*.]
		\item The absence of recent enforcement actions from the local regulator, which demonstrates that SunCoast is in satisfactory standing.
		\item SunCoast's servicing of unlicensed MSBs and provision of payable-through access to undisclosed sub-agents, combined with a four-year audit gap. These factors collectively mean NorthHarbor would have no ability to assess the AML quality of SunCoast's highest-risk customer segments, and no recent independent verification that SunCoast's controls are effective.
		\item The moderate FATF rating of SunCoast's home jurisdiction, which is a public information factor already priced into NorthHarbor's risk model.
		\item SunCoast's servicing of 14 MSBs, which is within the acceptable industry range for a commercial bank of its asset size.
		\item The questionnaire response itself, because submitting a CBDDQ indicates SunCoast's willingness to be transparent and therefore reduces rather than increases the risk.
	\end{enumerate}
	\textbf{Correct Answer: B}
	
	%--- Q14 ---
	\subsection{14. UNSCR 1540 and Proliferation Financing in Trade Finance \textnormal{[SELECT ONE]}}
	
	A trade finance bank in Singapore is processing a letter of credit for a transaction involving the export of 12 industrial vacuum pumps, 400 meters of maraging steel tubing, and a precision coordinate measuring machine. The declared end-user is a private engineering research institute in Country K, which is subject to active UN Security Council sanctions for its nuclear weapons development program. The exporter is a German industrial manufacturer and the stated end-use on the export license application is ``quality control laboratory equipment for automotive manufacturing.'' The Singapore bank's dual-use goods screening system flags two of the three items as controlled under EU and Wassenaar Arrangement export control lists.
	
	What is the most appropriate regulatory classification of this transaction, and what is the bank's primary obligation?
	
	\begin{enumerate}[label=\Alph*.]
		\item The transaction is a standard trade finance risk and requires only enhanced customs documentation before the letter of credit is issued.
		\item The transaction presents proliferation financing red flags. Maraging steel tubing, vacuum pumps, and coordinate measuring machines are recognized dual-use components in uranium enrichment and nuclear weapons manufacturing supply chains. The declared automotive end-use is inconsistent with the technical specifications of the flagged items. The bank must decline to issue the letter of credit, file an STR with Singapore's FIU (the Suspicious Transaction Reporting Office), and report to the relevant trade control authority under UNSCR 1540.
		\item The bank's only obligation is to verify that the German exporter holds a valid EU export license and confirm that the stated end-user is registered in Country K's corporate registry.
		\item Proliferation financing obligations apply only to financial institutions in the weapons-originating country (Germany); Singapore's bank has no independent obligation.
		\item The bank should process the transaction because Singapore is not a party to the Wassenaar Arrangement and is therefore not bound by its export control lists.
	\end{enumerate}
	\textbf{Correct Answer: B}
	
	%==============================================================
	\section{Part 3: AML Compliance Program Design and Risk-Based Approach}
	%==============================================================
	
	%--- Q15 ---
	\subsection{15. Enterprise-Wide Risk Assessment: Residual Risk Determination \textnormal{[SELECT ONE]}}
	
	A mid-sized regional bank is preparing its annual Enterprise-Wide Risk Assessment (EWRA). The assessment team calculates the following inherent risk scores for its three primary business lines:
	
	\begin{itemize}
		\item Retail banking: inherent risk score 6.2/10 (high cash transaction volume, significant MSB customer base)
		\item Commercial lending: inherent risk score 4.1/10 (standard corporate borrowers, moderate geographic risk)
		\item Private banking: inherent risk score 8.7/10 (international HNWIs, multiple PEP accounts, complex trust structures)
	\end{itemize}
	
	After applying mitigating controls (enhanced due diligence procedures, automated transaction monitoring, dedicated private banking compliance team), the residual risk scores are: retail 4.4/10; commercial lending 2.8/10; private banking 6.9/10. The EWRA methodology defines residual risk above 6.0 as ``high residual risk'' requiring escalation to the board and independent audit prioritization.
	
	Which of the following correctly describes the bank's next required steps?
	
	\begin{enumerate}[label=\Alph*.]
		\item The EWRA is complete once residual scores are documented. No further action is required as long as the bank's controls remain in place.
		\item The retail and private banking business lines both require board-level escalation. Retail's residual score of 4.4/10 exceeds the 4.0 threshold for medium-high risk, and private banking's 6.9/10 exceeds the 6.0 high risk threshold.
		\item Only the private banking business line requires board-level escalation, because it is the only line with a residual risk score exceeding the 6.0 high-risk threshold. The board must be informed, independent audit must prioritize this business line, and the compliance team must develop and implement a risk mitigation plan within a defined timeframe.
		\item All three business lines require immediate account closure of all high-risk relationships, because the EWRA has identified elevated risk across the institution.
		\item Commercial lending requires the highest priority attention because its inherent risk score declined the most from 4.1 to 2.8, suggesting the controls may be masking underlying suspicious activity.
	\end{enumerate}
	\textbf{Correct Answer: C}
	
	%--- Q16 ---
	\subsection{16. Three Lines of Defense: Independence Violation \textnormal{[SELECT ONE]}}
	
	A bank's internal audit team (Third Line of Defense) is asked by the Chief Compliance Officer to assist in redesigning the bank's AML transaction monitoring scenario library, because the compliance team (Second Line) lacks the technical expertise to build new detection scenarios. The internal audit team agrees and spends three months co-designing 14 new monitoring scenarios. Six months later, the same audit team conducts an independent test of the transaction monitoring system and issues a report concluding that all 14 new scenarios are effectively calibrated.
	
	What fundamental governance failure has occurred?
	
	\begin{enumerate}[label=\Alph*.]
		\item The Chief Compliance Officer violated GDPR by sharing transaction monitoring data with the internal audit team without customer consent.
		\item The internal audit team has compromised its independence by designing the same controls it subsequently audited. An audit team cannot objectively test controls it helped create. This violates the Three Lines of Defense framework and renders the audit findings unreliable, as the team cannot be objective about the adequacy of its own design decisions.
		\item The compliance team (Second Line) acted appropriately by delegating technical design work to internal audit; this is standard industry practice for small compliance teams.
		\item The failure is purely a resource management issue, not a governance failure, because the audit team still performed the audit as scheduled.
		\item The governance violation can be cured by having the Chief Compliance Officer sign an attestation confirming the audit findings are independent.
	\end{enumerate}
	\textbf{Correct Answer: B}
	
	%--- Q17 ---
	\subsection{17. Transaction Monitoring Calibration: False Negatives and Threshold Management \textnormal{[SELECT TWO]}}
	
	An AML operations director at a large retail bank is reviewing the results of a transaction monitoring system validation exercise. The validation team tested 500 historical SAR cases filed over a three-year period and found that 387 (77\%) of these cases would not have generated any system alert under the current monitoring thresholds and scenario parameters. The director argues that this finding is acceptable because the bank's alert volumes are already high (22,000 alerts per month), and raising sensitivity would create an unmanageable operational burden. A new external audit, however, rates the monitoring system as ``ineffective'' in three key scenario categories.
	
	Which two of the following assessments are correct regarding this situation?
	
	\begin{enumerate}[label=\Alph*.]
		\item A 77\% false-negative rate on known SAR cases demonstrates that the current calibration is systematically missing genuine suspicious activity, which constitutes a material deficiency in the AML program that exposes the institution to significant regulatory risk.
		\item Operational manageability (the volume of alerts) is a valid independent justification for maintaining threshold levels that produce high false-negative rates, even if regulators disagree.
		\item The monitoring system must be recalibrated to reduce the false-negative rate. The bank should also document the validation findings, the identified deficiencies, a remediation plan with timelines, and present the plan to senior management and the board.
		\item A monitoring system that generates 22,000 alerts per month is inherently effective, regardless of false-negative rates, because high alert volume demonstrates comprehensive surveillance.
		\item The external audit's rating of ``ineffective'' on three scenario categories can be disregarded if the bank can demonstrate that the three categories relate to low-dollar-value transactions.
	\end{enumerate}
	\textbf{Correct Answers: A and C}
	
	%--- Q18 ---
	\subsection{18. PEP Lifecycle Management: Mid-Relationship Designation Change \textnormal{[SELECT ONE]}}
	
	A commercial bank has maintained an account for the past eight years for Mr. T, a successful entrepreneur with a domestic manufacturing business. The account has been classified as medium risk since opening. During a routine KYC refresh, a compliance analyst discovers through an OSINT review that Mr. T was appointed six months ago as the Deputy Minister of Finance in a foreign government. The appointment was publicly announced in the foreign government's official gazette. Mr. T did not proactively disclose this change to the bank, and the account's risk rating has not been updated since the appointment.
	
	Which of the following best describes the correct compliance response and any violations that have occurred?
	
	\begin{enumerate}[label=\Alph*.]
		\item Because Mr. T did not proactively disclose the change, the bank has no obligation to update his risk profile until Mr. T personally provides notification in writing.
		\item The compliance analyst's discovery triggers immediate obligations. Mr. T must be reclassified as a Foreign PEP, the account must be elevated to high risk with Enhanced Due Diligence, the source of wealth and source of funds must be re-established and verified, and senior management approval must be obtained to continue the relationship. The six-month gap since the appointment also represents a control failure in the bank's ongoing monitoring and KYC refresh program, which should be reviewed.
		\item Because Mr. T's role is Deputy Minister rather than Minister, he does not qualify as a PEP under FATF standards, which only cover senior government officials.
		\item The bank's obligation is limited to sending Mr. T a letter requesting that he update his customer profile at his earliest convenience.
		\item Because the account has been clean for eight years and the manufacturing business is established, no change in risk classification is warranted.
	\end{enumerate}
	\textbf{Correct Answer: B}
	
	%--- Q19 ---
	\subsection{19. Beneficial Ownership: Aggregating Direct and Indirect Stakes \textnormal{[SELECT ONE]}}
	
	A compliance officer is onboarding a new client, Apex Industrial Corp., applying for a trade finance facility. Apex's ownership structure is: Individual K owns 18\% of Apex directly. Individual K also owns 90\% of Zenith Holdings LLC, which in turn owns 15\% of Apex. The regulatory UBO threshold is 25\% combined direct and indirect ownership.
	
	What is Individual K's total effective ownership stake in Apex, and what is the correct compliance treatment?
	
	\begin{enumerate}[label=\Alph*.]
		\item Individual K's effective ownership is 18\% (direct). The 15\% held through Zenith Holdings cannot be attributed to Individual K without a court order confirming common control. Total: 18\%. Individual K is below the 25\% threshold and does not need to be verified as a UBO.
		\item Individual K's effective ownership is calculated as: 18\% direct plus (90\% $\times$ 15\%) = 18\% + 13.5\% = 31.5\%. Because 31.5\% exceeds the 25\% threshold, Individual K must be identified, verified, and documented as a UBO of Apex Industrial Corp.
		\item Individual K's effective ownership is 33\% (18\% + 15\% pass-through), which exceeds the 25\% threshold and requires UBO verification.
		\item Individual K does not qualify as a UBO because the indirect stake is held through a separate legal entity (Zenith Holdings) that must itself be treated as the intermediate beneficial owner.
		\item The 25\% threshold applies only to direct ownership; indirect stakes held through intermediate entities are irrelevant to the UBO calculation under FinCEN's CDD Final Rule.
	\end{enumerate}
	\textbf{Correct Answer: B}
	
	%--- Q20 ---
	\subsection{20. AML Training: Role-Based Requirements \textnormal{[SELECT TWO]}}
	
	A bank's compliance officer is redesigning the annual AML training curriculum. She must assign training modules to four distinct employee groups: (1) frontline tellers, (2) private banking relationship managers, (3) the board of directors, and (4) the IT infrastructure team responsible for maintaining the transaction monitoring software servers. The prior year's training delivered the same two-hour online module to all employees and was rated ``ineffective'' by the bank's external auditor because training lacked relevance to specific job functions.
	
	Which two of the following training design principles are most critical for regulatory compliance under CAMS standards?
	
	\begin{enumerate}[label=\Alph*.]
		\item Training must be tailored to the specific roles, responsibilities, and ML/TF risks faced by each employee group. Frontline tellers need training on cash structuring indicators and CTR reporting; private bankers need PEP, source of wealth, and EDD training; board members need governance oversight and ML/TF risk strategy training.
		\item All employees, regardless of role, must complete an identical training program so that the bank can demonstrate consistent institutional coverage and cannot be accused of differential training standards.
		\item The IT infrastructure team maintaining server hardware has no material exposure to AML risk and can be exempted from all AML training without regulatory consequence.
		\item Training frequency and content must be documented and presented to the board on at least an annual basis, with evidence that training effectiveness has been assessed.
		\item AML training is only required for employees who have direct customer contact; back-office and technology staff have no training obligations.
	\end{enumerate}
	\textbf{Correct Answers: A and D}
	
	%==============================================================
	\section{Part 4: Investigations, SARs, and Law Enforcement Cooperation}
	%==============================================================
	
	%--- Q21 ---
	\subsection{21. SAR Filing Triggers: The Alert Lifecycle and Investigative Standards \textnormal{[SELECT ONE]}}
	
	A transaction monitoring analyst at a U.S. bank reviews an automated alert flagging 23 cash deposits into a business account over 45 days, each between \$8,400 and \$9,700. The account belongs to a licensed vehicle dealership. The analyst notes in the system: ``Customer is a car dealer. Cash from car sales is normal. No SAR needed.'' and closes the alert without requesting any supporting documentation from the branch, reviewing the account's expected activity profile, checking whether the deposit pattern is consistent with known customer behavior, or consulting with the AML supervisor. No additional review is performed.
	
	What is the most significant AML compliance failure in the analyst's approach, and what is the correct course of action?
	
	\begin{enumerate}[label=\Alph*.]
		\item The analyst correctly applied the risk-based approach by using contextual business knowledge to close the alert. No failure occurred.
		\item The analyst committed the fundamental error of closing an alert based on unverified assumptions. The deposit pattern is a textbook structuring red flag and the automobile dealership context does not automatically legitimize repeated sub-threshold cash deposits. The analyst must obtain and review source documents (daily sales logs, title transfer records, customer receipts), compare deposits against expected cash flow for a dealership of this size, and escalate to a supervisor if a legitimate explanation cannot be documented. Failure to complete this analysis before closing the alert is a material AML program deficiency.
		\item The failure is procedural only: the analyst should have used a different dropdown code in the case management system to close the alert.
		\item The analyst should have immediately filed a SAR based solely on the pattern of deposits without conducting any further investigation, because structuring is evident on its face.
		\item The failure is that the analyst should have filed a CTR for each deposit rather than a SAR, because the transactions involve cash.
	\end{enumerate}
	\textbf{Correct Answer: B}
	
	%--- Q22 ---
	\subsection{22. Tipping-Off: Permissible vs.\ Prohibited Disclosures \textnormal{[SELECT TWO]}}
	
	A bank's compliance team has filed an STR on Mr. V, a wealthy importer, based on a pattern of large cash deposits inconsistent with his stated import business. Three scenarios occur simultaneously: (i) Mr. V calls his relationship manager asking why a wire transfer has been delayed; the relationship manager tells him the compliance team is ``reviewing some transactions'' on his account; (ii) the compliance team shares the STR narrative with its legal department as part of an internal legal review; (iii) a second bank where Mr. V also holds an account contacts the compliance department as part of a joint investigation into suspected organized crime; the compliance team shares its STR findings under a legally permissible information-sharing arrangement.
	
	Which two of the following accurately identify a tipping-off violation?
	
	\begin{enumerate}[label=\Alph*.]
		\item Scenario (i) constitutes tipping-off. The relationship manager's statement that the compliance team is ``reviewing transactions'' alerts Mr. V that he is under investigation, which is a violation of AML laws regardless of whether a formal STR has been filed.
		\item Scenario (ii) constitutes tipping-off because disclosing the STR to the bank's own legal department spreads knowledge of the filing beyond the filing officer.
		\item Scenario (iii) constitutes tipping-off because sharing STR information with another bank, even in an inter-institution AML investigation, is never permissible.
		\item All three scenarios constitute tipping-off violations because any disclosure related to an STR, for any purpose, is prohibited.
		\item Scenario (i) constitutes a tipping-off violation even though the relationship manager did not mention the STR by name, because disclosing the existence of compliance scrutiny of a specific account to the account holder is sufficient to constitute an impermissible disclosure.
	\end{enumerate}
	\textbf{Correct Answers: A and E}
	
	%--- Q23 ---
	\subsection{23. MLAT and FIU Cooperation: Egmont Group Principles \textnormal{[SELECT ONE]}}
	
	An FIU in Country Alpha has identified a multi-layered embezzlement scheme in which a state-owned enterprise's CFO siphoned \$47 million in public funds over six years. The funds were routed through a chain of shell companies in four different jurisdictions before being invested in private equity funds in Country Beta. Country Alpha and Country Beta have an active Mutual Legal Assistance Treaty (MLAT) and are both members of the Egmont Group. Country Alpha's FIU wants to obtain account records from financial institutions in Country Beta to support a criminal prosecution. However, the process must be handled through official channels because Country Beta's banking secrecy laws prevent voluntary disclosure by its banks.
	
	Which mechanism provides the most appropriate and legally binding pathway for Country Alpha to obtain the account records from Country Beta?
	
	\begin{enumerate}[label=\Alph*.]
		\item Country Alpha's FIU should contact Country Beta's FIU through the Egmont Group's secure communication platform (Egmont Secure Web) to exchange financial intelligence at the FIU-to-FIU level, while Country Alpha's law enforcement simultaneously submits a formal MLAT request to Country Beta's competent judicial authority for the production of bank account records that will be admissible in criminal proceedings.
		\item Country Alpha can instruct its domestic financial institutions to demand account records directly from Country Beta's banks using a private contractual arrangement, without involving any government authorities.
		\item Country Alpha's law enforcement agency should publicly name Country Beta's banks in a press conference to pressure them to voluntarily disclose the records.
		\item The Egmont Group has binding enforcement authority over all member FIUs and can compel Country Beta's FIU to produce the bank account records within 10 business days.
		\item Country Alpha must first obtain a conviction in its domestic courts before requesting any international cooperation.
	\end{enumerate}
	\textbf{Correct Answer: A}
	
	%--- Q24 ---
	\subsection{24. Asset Forfeiture: Civil vs.\ Criminal Standards \textnormal{[SELECT ONE]}}
	
	A federal prosecutor in the United States is pursuing a money laundering case against a criminal network that used a series of shell companies and real estate purchases to launder \$22 million in drug trafficking proceeds. The prosecutor is considering whether to proceed under criminal forfeiture (requiring conviction of the underlying defendant) or civil asset forfeiture (proceeding directly against the property). The defendants' attorneys argue that all of the properties are titled in the names of innocent third-party shell companies that were unaware of the laundering.
	
	Which of the following most accurately describes the key distinction between criminal and civil asset forfeiture in this context?
	
	\begin{enumerate}[label=\Alph*.]
		\item Criminal forfeiture requires conviction of the named defendant and is filed as part of the criminal prosecution. Civil forfeiture (in rem) proceeds against the property itself rather than against a person, requires the government to demonstrate by a preponderance of the evidence that the property was involved in or proceeds of criminal activity, and does not require a conviction. Civil forfeiture can proceed even if the defendant is acquitted or cannot be located.
		\item Criminal forfeiture requires a lower standard of proof than civil forfeiture and is therefore preferred by prosecutors in complex money laundering cases.
		\item Civil forfeiture is only available for drug trafficking proceeds and cannot be applied to money laundering proceeds derived from other predicate offenses.
		\item Civil forfeiture automatically transfers title of the seized property to the federal government without any judicial hearing, because the property is presumed guilty under U.S. law.
		\item Criminal forfeiture allows the government to seize property owned by innocent third parties without any notice or hearing, as long as a grand jury indictment has been issued.
	\end{enumerate}
	\textbf{Correct Answer: A}
	
	%--- Q25 ---
	\subsection{25. SAR Recordkeeping, Safe Harbor, and Confidentiality \textnormal{[SELECT TWO]}}
	
	A compliance officer at a U.S. bank is briefing the board of directors on SAR program obligations. She presents three statements for board ratification: (1) ``A filed SAR and all supporting documentation must be retained for five years from the date of filing''; (2) ``If a customer sues the bank for filing a SAR, the bank is immune from civil liability if the SAR was filed in good faith, but the bank may use the SAR as evidence to defend itself in court''; (3) ``A bank employee who discloses the existence of a SAR to the customer or any unauthorized third party is subject to criminal penalties.''
	
	Which two of these statements are accurate?
	
	\begin{enumerate}[label=\Alph*.]
		\item Statement (1) is accurate. Under 31 U.S.C.\ \S 5318(g)(3), SAR records and supporting documentation must be retained for five years from the date of the report.
		\item Statement (2) is accurate. The SAR safe harbor provides immunity from civil and criminal liability for good-faith filers. However, the second part of Statement (2) is \textit{inaccurate}: a SAR and its contents are protected by a strict confidentiality requirement and may not be disclosed or used in court proceedings by the filing institution.
		\item Statement (3) is accurate. Unauthorized disclosure of a SAR filing is a criminal offense under 31 U.S.C.\ \S 5318(g)(2), punishable by fines and imprisonment.
		\item Statement (1) is inaccurate; SAR records only need to be retained for three years under the BSA.
		\item Statement (3) is inaccurate; only intentional disclosure carries criminal liability; inadvertent disclosure is only subject to a civil fine.
	\end{enumerate}
	\textbf{Correct Answers: A and C}
	
	%==============================================================
	\section{Part 5: KYC, CDD, EDD, and Customer Risk}
	%==============================================================
	
	%--- Q26 ---
	\subsection{26. Customer Due Diligence: Beneficial Ownership of Trusts \textnormal{[SELECT ONE]}}
	
	A private bank is onboarding a discretionary family trust established under the laws of a British Crown Dependency. The trust's key parties are: a corporate trustee (a licensed trust company); a protector (Mr. F, a retired former head of state of an emerging market country); and a class of beneficiaries comprising five adult members of the same family. The trust corpus is \$28 million, consisting of foreign real estate and a diversified securities portfolio. The trust deed grants the trustee full discretionary power over distributions, but the protector has a veto right over all distributions above \$200,000 and the right to replace the trustee.
	
	Under CAMS and FATF standards, which parties must the bank identify and verify, and at what risk level?
	
	\begin{enumerate}[label=\Alph*.]
		\item Only the corporate trustee requires identification, because the trustee is the legal owner of the trust assets and bears all legal obligations as the party opening the account.
		\item The bank must identify and verify: the corporate trustee and its UBOs; the protector (Mr. F) who, as a former head of state, qualifies as a Politically Exposed Person requiring Enhanced Due Diligence and senior management approval; all five beneficiaries to the extent they are determinable; and the settlor who established the trust. The entire relationship must be classified as high risk given the PEP protector status, the offshore trust jurisdiction, and the significant asset value.
		\item Only the beneficiaries need to be identified, because they are the parties with the ultimate economic interest in the trust.
		\item The trust qualifies for Simplified Due Diligence because it is established under the laws of a Crown Dependency, which is a well-regulated financial jurisdiction.
		\item The protector does not require identification because his veto right is a passive governance mechanism and does not constitute ownership or control of the trust.
	\end{enumerate}
	\textbf{Correct Answer: B}
	
	%--- Q27 ---
	\subsection{27. Enhanced Due Diligence: Source of Wealth vs.\ Source of Funds \textnormal{[SELECT ONE]}}
	
	A private banker is conducting Enhanced Due Diligence on a PEP client, the Minister of Natural Resources of a West African country. The minister wishes to deposit \$4.5 million, claiming the funds represent the sale proceeds of a residential property in Paris. He provides a notarized copy of the sale contract showing \$4.5 million was received. When asked about source of wealth (how he accumulated his overall wealth), the minister states he has been a senior civil servant all his adult life and also inherited property from his parents. He declines to provide documentation of his net worth or of the parental inheritance. The private banker proposes to accept the source of funds documentation (the property sale contract) as sufficient and open the account.
	
	Why is the banker's proposed approach insufficient under EDD standards?
	
	\begin{enumerate}[label=\Alph*.]
		\item The property sale contract is not an acceptable source of funds document; only bank statements qualify.
		\item EDD for PEPs requires establishing both the source of funds (the specific transaction being deposited, i.e., the property sale) and the source of wealth (how the PEP accumulated all of their assets overall). The sale contract addresses only source of funds. The critical unresolved question -- how a lifelong civil servant on a public salary accumulated a \$4.5 million Paris property and sufficient other wealth to generate this kind of asset base -- must be answered through corroborative documentation. The minister's refusal to provide supporting documentation is itself a red flag that must be escalated to senior management before the account can be opened.
		\item The minister's refusal to provide documentation is a legal right and cannot be used as a red flag in the EDD process.
		\item The banker should simply add a notation in the system that the minister declined to provide wealth documentation, downgrade the overall risk rating to reflect the clean source of funds, and proceed.
		\item EDD for a government minister is only required if the minister's country is on the FATF Grey List.
	\end{enumerate}
	\textbf{Correct Answer: B}
	
	%--- Q28 ---
	\subsection{28. Ongoing Monitoring: Triggering a KYC Refresh \textnormal{[SELECT ONE]}}
	
	A commercial bank's CDD policy requires a full KYC refresh for medium-risk corporate customers every three years. A transaction monitoring alert is generated for the account of Bravo Logistics LLC, a domestic freight company classified as medium risk at onboarding two years ago. The alert identifies 14 international wire transfers to counterparties in Syria and North Korea, neither of which appeared in the company's original customer profile. When the branch relationship manager contacts the company, she is told by an unfamiliar employee that ``management has changed and the business is now doing some international work.'' The relationship manager accepts this explanation and closes the alert with a note: ``customer confirmed legitimate international business; next scheduled KYC refresh in 12 months.''
	
	What is the relationship manager's critical compliance error?
	
	\begin{enumerate}[label=\Alph*.]
		\item The relationship manager applied the risk-based approach correctly by documenting the customer's explanation and scheduling the next KYC refresh within the standard policy timeframe.
		\item The relationship manager failed to trigger an unscheduled, immediate KYC refresh. Material changes in a customer's business profile -- specifically the introduction of wire transfers to OFAC and UN-sanctioned jurisdictions (Syria and North Korea), a change in management, and new international activity -- constitute trigger events requiring immediate out-of-cycle CDD review and EDD, regardless of the scheduled three-year cycle. The transactions involving Syria and North Korea must also be escalated immediately for sanctions screening review and potential SAR filing.
		\item The relationship manager's only error was not escalating the change of management to the branch manager; the international business change is not a CDD trigger event.
		\item The correct action was to immediately close the account, because wire transfers to sanctioned countries are automatically disqualifying.
		\item The relationship manager should have filed a CTR for each wire transfer rather than treating it as a KYC issue.
	\end{enumerate}
	\textbf{Correct Answer: B}
	
	%--- Q29 ---
	\subsection{29. Simplified Due Diligence: When It Cannot Apply \textnormal{[SELECT ONE]}}
	
	A compliance officer is reviewing the CDD classification of five account relationships to determine whether Simplified Due Diligence (SDD) is appropriate for any of them:
	
	\begin{itemize}
		\item Entity 1: A domestic municipality opening a public payroll account.
		\item Entity 2: A subsidiary of a publicly listed company, but the subsidiary itself is incorporated in a FATF Black List jurisdiction and its operations are opaque.
		\item Entity 3: A domestic pension fund supervised by the national pension regulator.
		\item Entity 4: A domestic credit union supervised by the national financial regulator.
		\item Entity 5: A foreign government-owned enterprise operating in an oil sector with high corruption risk.
	\end{itemize}
	
	For which entities is SDD clearly \textit{inappropriate} under a risk-based AML framework?
	
	\begin{enumerate}[label=\Alph*.]
		\item Entities 1, 3, and 4 are inappropriate for SDD because all government-linked entities carry elevated corruption risk.
		\item Entities 2 and 5 are inappropriate for SDD. Entity 2 cannot benefit from its parent's public listing when it itself is incorporated in a FATF Black List jurisdiction with opaque operations -- the regulatory transparency that justifies SDD does not extend to opaque subsidiaries in non-compliant jurisdictions. Entity 5's combination of government ownership, oil sector exposure, and high corruption risk requires EDD, not SDD.
		\item Only Entity 2 is inappropriate for SDD; Entity 5 qualifies for SDD because it is government-owned and therefore publicly accountable.
		\item None of the five entities are appropriate for SDD because SDD is only available in exceptional circumstances approved by regulators in advance.
		\item All five entities qualify for SDD because they are either government-related or supervised by a regulator.
	\end{enumerate}
	\textbf{Correct Answer: B}
	
	%--- Q30 ---
	\subsection{30. Sanctions Screening: Fuzzy Matching and Name Variant Obligations \textnormal{[SELECT TWO]}}
	
	A compliance analyst at a U.S. payment processor is reviewing a wire transfer initiated by a sender whose name is listed as ``Muammar Gadaffi'' in the transaction data. The OFAC SDN list contains the name ``Moammar Qaddafi'' as the primary listing, along with 48 alternative name transliterations in the entity's SDN listing. The payment processor's sanctions screening system uses exact-string matching and did not generate an alert because no character sequence matched precisely. The wire transfer was processed. An internal review identifies the missed match two weeks later.
	
	Which two of the following conclusions are correct?
	
	\begin{enumerate}[label=\Alph*.]
		\item The payment processor's use of exact-string matching for Arabic-origin names with many common transliterations is an inadequate screening methodology. Sanctions screening systems must employ fuzzy logic or algorithmic name-matching capable of identifying phonetically similar and transliterated name variants to meet OFAC compliance standards.
		\item The payment processor bears no liability because OFAC's SDN entry did not exactly match the wire transfer data, and exact-string matching is a legally defensible standard under the BSA.
		\item The missed match constitutes a potential OFAC sanctions violation. The payment processor must file a Voluntary Self-Disclosure with OFAC, block or reject the transaction retroactively (to the extent possible), and implement corrective action to its screening methodology.
		\item The analyst should simply file a SAR and close the matter without any obligation to report to OFAC.
		\item The payment processor is not subject to OFAC obligations because it processes payments rather than initiating them; the originating bank bears sole responsibility.
	\end{enumerate}
	\textbf{Correct Answers: A and C}
	
	%==============================================================
	\section{Part 6: Virtual Assets, DeFi, and Emerging Typologies}
	%==============================================================
	
	%--- Q31 ---
	\subsection{31. FATF Travel Rule: VASP-to-VASP Transfers \textnormal{[SELECT ONE]}}
	
	CryptoVault, a registered VASP, receives a transfer of 4.2 Bitcoin (current value: approximately \$180,000) from ExchangeX, another VASP. The Bitcoin arrives with no originator or beneficiary information attached -- no name, wallet address attribution, or account reference. CryptoVault's compliance system flags the transfer as a Travel Rule data deficiency. The receiving customer states that the funds are a personal Bitcoin investment. CryptoVault's compliance officer argues that because Bitcoin is a bearer asset and the customer has already passed KYC at CryptoVault, the missing originator data from ExchangeX is an ExchangeX compliance problem, not CryptoVault's.
	
	Why is this analysis incorrect under FATF's Virtual Assets guidance?
	
	\begin{enumerate}[label=\Alph*.]
		\item The analysis is correct: the Travel Rule imposes obligations only on the originating VASP, not the receiving VASP.
		\item The analysis is incorrect. Under the FATF Travel Rule (Recommendation 16 as updated for virtual assets), the receiving VASP has an obligation to obtain and hold the required originator and beneficiary information. If the information is not provided by the sending VASP, the receiving VASP must take follow-up action, which may include requesting the information from the sending VASP, delaying the transaction, and applying enhanced scrutiny or SAR filing if the information cannot be obtained. CryptoVault cannot simply accept the transfer as compliant because its own KYC on the beneficiary does not substitute for originator data from the sending institution.
		\item The Travel Rule applies only to fiat-to-crypto or crypto-to-fiat transactions; Bitcoin-to-Bitcoin transfers between two VASPs are exempt.
		\item The FATF Travel Rule threshold for virtual assets is \$1 million per transaction; the \$180,000 transfer is below this threshold and therefore Travel Rule obligations do not apply.
		\item Because CryptoVault is the receiving VASP, its only obligation is to verify the identity of the beneficiary, which it has already done through KYC. Originator data is exclusively the sending VASP's responsibility.
	\end{enumerate}
	\textbf{Correct Answer: B}
	
	%--- Q32 ---
	\subsection{32. DeFi and AML Responsibility: Beneficial Ownership of a Protocol \textnormal{[SELECT ONE]}}
	
	A group of software developers created a decentralized lending protocol (DLP) two years ago. The protocol is governed by a DAO (Decentralized Autonomous Organization) in which governance token holders vote on protocol changes. The five founding developers retain 68\% of governance tokens and continue to upgrade the protocol's code. The DLP has processed over \$2 billion in transactions. No licensing or registration with any financial regulator has been sought. Law enforcement in Country X has identified the DLP as the primary layering mechanism in a \$90 million drug trafficking money laundering case.
	
	Under FATF guidance on virtual assets, which statement most accurately describes the AML obligations of the founding developers?
	
	\begin{enumerate}[label=\Alph*.]
		\item The founding developers have no AML obligations because the protocol operates autonomously through smart contracts and no human intermediary processes the transactions.
		\item Under FATF guidance, where a person or entity maintains control or sufficient influence over a DeFi protocol -- including through governance token majority, the ability to upgrade contract code, or ongoing management -- they may be considered a VASP and subject to AML/CFT obligations. The founding developers' retention of 68\% of governance tokens and continued code control means they exercise significant influence and are likely subject to registration, KYC, transaction monitoring, and SAR filing obligations.
		\item Only the DAO token holders who voted in favor of the money laundering-related protocol changes face AML liability; the founders are insulated by the DAO structure.
		\item The DLP's AML obligations rest entirely with the borrowers and lenders who use the protocol, not with its developers or governance structure.
		\item FATF guidance on DeFi is purely advisory; no country has enacted binding legislation applying VASP obligations to DeFi protocol developers.
	\end{enumerate}
	\textbf{Correct Answer: B}
	
	%--- Q33 ---
	\subsection{33. NFT Market Abuse: Wash Trading Identification \textnormal{[SELECT TWO]}}
	
	A blockchain analytics firm provides a report to a regulated NFT marketplace showing the following activity: Wallet Group A (comprising 17 wallets) has collectively purchased and resold a series of 12 NFTs among themselves 89 times over 60 days, with sale prices escalating from an initial \$2,000 per NFT to a peak of \$340,000 per NFT. Blockchain analysis confirms that all 17 wallets are funded from the same source wallet at a centralized exchange, where the KYC-verified account holder is Individual P. Individual P then sells one NFT to an unconnected third-party buyer for \$310,000 in ETH and converts the ETH to fiat currency through the centralized exchange, reporting the proceeds as a capital gain from ``investment art.''
	
	Which two of the following are accurate findings?
	
	\begin{enumerate}[label=\Alph*.]
		\item The activity constitutes NFT wash trading, where self-controlled wallets artificially inflate the price of an asset through circular transactions to create a fraudulent market record, which is then used to launder illicit funds by selling the NFT to an unwitting third party at the artificially inflated price.
		\item The centralized exchange that funded all 17 wallets from a single KYC-verified account should file a SAR based on the wash trading pattern, because the circular transactions are a layering indicator regardless of whether the underlying proceeds are from a predicate crime or from legitimate funds being manipulated for market fraud.
		\item Wash trading through NFTs is not a recognized money laundering typology and does not trigger SAR filing obligations.
		\item Because all wallets belonged to the same verified individual, the transactions are intra-account transfers with no third-party AML risk.
		\item Individual P's conversion of \$310,000 to fiat and reporting it as a capital gain fully satisfies all AML obligations because the transaction is on the tax record.
	\end{enumerate}
	\textbf{Correct Answers: A and B}
	
	%--- Q34 ---
	\subsection{34. Online Gaming and Virtual Currency: Placement Through In-Game Economies \textnormal{[SELECT ONE]}}
	
	A law enforcement financial intelligence analyst is investigating a suspected money laundering network that has exploited an online multiplayer game's in-game currency system. The scheme works as follows: criminal proceeds (approximately \$600,000 in cash) are used to purchase in-game currency through third-party ``gold selling'' websites using prepaid debit cards. The in-game currency is then distributed across 40 different gamer accounts, each of which sells the currency back to other players through a player-to-player marketplace. Each sale is received as a small PayPal payment (\$200--\$800 per transaction). The PayPal proceeds are deposited into clean bank accounts and reported as ``freelance gaming income.''
	
	At which stage of the money laundering cycle does the purchase of in-game currency using prepaid debit cards most accurately occur?
	
	\begin{enumerate}[label=\Alph*.]
		\item Integration, because the in-game currency has already been separated from the criminal proceeds by the time it is purchased.
		\item Placement, because the criminal cash proceeds are being introduced into a quasi-financial system (the in-game economy) by purchasing in-game currency. The use of prepaid cards adds a layer of anonymity to the initial placement.
		\item Layering, because the use of prepaid cards adds an intermediate step between the cash and the in-game currency.
		\item Structuring, because the PayPal payments from game sales are consistently kept below \$1,000.
		\item Predicate offense, because operating a gold-selling website may violate the game publisher's terms of service.
	\end{enumerate}
	\textbf{Correct Answer: B}
	
	%==============================================================
	\section{Part 7: Sanctions, OFAC, and Proliferation Financing}
	%==============================================================
	
	%--- Q35 ---
	\subsection{35. OFAC 50\% Rule: Multi-Tier Ownership Calculation \textnormal{[SELECT ONE]}}
	
	A bank's sanctions compliance team is analyzing the ownership structure of GammaCo, a target entity in a pending wire transfer. The ownership chain is as follows:
	
	\begin{itemize}
		\item Individual X (OFAC SDN-listed) owns 55\% of AlphaCorp.
		\item AlphaCorp owns 40\% of BetaGroup.
		\item BetaGroup owns 60\% of GammaCo.
		\item Individual Y (non-sanctioned) owns 40\% of GammaCo directly.
	\end{itemize}
	
	Under OFAC's 50\% Rule, is GammaCo a blocked entity?
	
	\begin{enumerate}[label=\Alph*.]
		\item GammaCo is not blocked because Individual X does not directly own any interest in GammaCo.
		\item AlphaCorp is blocked because Individual X owns 55\%. BetaGroup is \textit{not} blocked because AlphaCorp only owns 40\% of BetaGroup (below 50\%). Because BetaGroup is not blocked, its 60\% ownership of GammaCo does not transmit SDN-blocked status to GammaCo. GammaCo is not a blocked entity under the OFAC 50\% Rule.
		\item AlphaCorp is blocked (55\% owned by SDN-listed Individual X). BetaGroup is blocked because AlphaCorp, a blocked entity, owns 40\% of BetaGroup -- and under OFAC's interpretation, any ownership by a blocked entity transmits blocked status to the owned entity regardless of the percentage.
		\item GammaCo is automatically blocked because the ownership chain traces back to an SDN-listed individual through any number of tiers.
		\item GammaCo is blocked because BetaGroup's 60\% ownership exceeds the 50\% threshold, and BetaGroup must be analyzed as a blocked entity because it is owned by AlphaCorp, which is blocked.
	\end{enumerate}
	\textbf{Correct Answer: B}
	
	%--- Q36 ---
	\subsection{36. Secondary Sanctions and Extra-Territorial Jurisdiction \textnormal{[SELECT TWO]}}
	
	A German bank, EuroCorp, processes a EUR-denominated loan of euro 45 million to a Turkish conglomerate. The loan proceeds are used to fund the construction of a natural gas extraction facility in Iran. The construction contract is held by an Iranian state-owned energy company not individually listed on the OFAC SDN list. No U.S. dollars are used, no U.S. persons are involved, and no transaction touches the U.S. financial system. EuroCorp's legal team argues that U.S. sanctions cannot apply because there is no U.S. nexus.
	
	Which two of the following statements correctly assess EuroCorp's risk?
	
	\begin{enumerate}[label=\Alph*.]
		\item EuroCorp's legal team is correct: OFAC primary sanctions only apply to U.S. persons and U.S.-dollar transactions. With no U.S. nexus, EuroCorp bears no OFAC exposure.
		\item EuroCorp faces U.S. secondary sanctions risk under the Iran Threat Reduction and Syria Human Rights Act (ITRSHRA) and related legislation. U.S. secondary sanctions can target non-U.S. entities that conduct ``significant'' transactions with designated Iranian energy sector entities, even without a U.S. dollar nexus, by restricting access to U.S. markets, the U.S. financial system, and correspondent banking relationships with U.S. institutions.
		\item EuroCorp is protected from secondary sanctions because the EU Blocking Statute prohibits EU companies from complying with U.S. secondary sanctions against Iran, creating a legal defense in EU courts.
		\item The practical risk for EuroCorp is that even if the EU Blocking Statute provides a defense in EU courts, U.S. regulators can still impose secondary sanctions that would restrict EuroCorp's USD correspondent banking access, effectively cutting it off from dollar-clearing markets and creating severe operational consequences regardless of EU law protections.
		\item EuroCorp's exposure is limited to EU sanctions on Iran, which only prohibit arms and nuclear material-related transactions; the natural gas sector transaction is permitted under current EU Iran sanctions.
	\end{enumerate}
	\textbf{Correct Answers: B and D}
	
	%--- Q37 ---
	\subsection{37. Dual-Use Goods and the Proliferation Financing Supply Chain \textnormal{[SELECT ONE]}}
	
	A trade finance compliance analyst at an Asian bank reviews a transaction involving the export of the following items by an electronics distributor: (a) 800 units of high-frequency oscilloscopes capable of measuring signals above 1 GHz; (b) 200 units of radiation hardened integrated circuits (RadHard ICs); (c) 500 meters of specialized fiber optic cable with vibration-dampening jacketing. The buyer is a ``scientific research institute'' in Country Z, which is under a UN Security Council arms embargo and active FATF monitoring for proliferation financing. The payment is routed through four intermediate companies across three jurisdictions. The exporter states the end-use is ``atmospheric research instrumentation.''
	
	Which factor most strongly indicates that this transaction constitutes a proliferation financing red flag?
	
	\begin{enumerate}[label=\Alph*.]
		\item The volume of items (800 oscilloscopes) suggests a commercial rather than a research purchase and may indicate resale, which is a trade discrepancy but not a proliferation indicator.
		\item The combination of RadHard ICs (which are controlled under Wassenaar Arrangement and MTCR export control lists due to their application in missile guidance, nuclear weapons detonation systems, and satellite platforms), the destination country's active UN arms embargo and FATF monitoring for proliferation financing, the multi-intermediary payment routing, and an end-use claim (atmospheric research) that is technically inconsistent with RadHard IC applications collectively constitutes a high-probability proliferation financing red flag requiring the bank to decline financing, file an STR, and notify the relevant export control authority.
		\item Only the fiber optic cable requires review; oscilloscopes and RadHard ICs are commercially available consumer electronics with no export control implications.
		\item The bank should process the transaction because the exporter holds a valid domestic export license; the responsibility for proliferation screening rests entirely with the exporting country's customs agency.
		\item Multi-intermediary payment routing is standard practice in international trade finance and does not constitute an independent red flag.
	\end{enumerate}
	\textbf{Correct Answer: B}
	
	%==============================================================
	\section{Part 8: Terrorist Financing and NPO Risk}
	%==============================================================
	
	%--- Q38 ---
	\subsection{38. Terrorist Financing: Key Distinctions from Money Laundering \textnormal{[SELECT ONE]}}
	
	A financial intelligence analyst is preparing a briefing for a bank's board on the key operational differences between money laundering (ML) and terrorist financing (TF). The analyst presents five statements and asks the board to identify which most accurately describes a critical distinction between the two financial crime types.
	
	\begin{enumerate}[label=\Alph*.]
		\item Money laundering always involves large sums of money, while terrorist financing always involves small sums.
		\item In money laundering, the criminal objective is to conceal the illicit origin of the proceeds and make them appear legitimate. In terrorist financing, the funds may originate from entirely legal sources (charitable donations, business income, personal savings) and the concealment objective is directed at the intended use (financing a terrorist act) rather than the origin. This distinction is critical because standard AML controls focused on suspicious origins may fail to detect TF involving clean-source funds.
		\item Terrorist financing is only relevant for banks with operations in conflict zones; domestic retail banks have no TF exposure.
		\item Both ML and TF involve the same three-stage cycle (placement, layering, integration); the only difference is the geographic location of the criminal network.
		\item Terrorist financing always involves wire transfers above \$10,000; money laundering typically involves structured cash deposits below the CTR threshold.
	\end{enumerate}
	\textbf{Correct Answer: B}
	
	%--- Q39 ---
	\subsection{39. NPO Risk: Distinguishing Legitimate Operations from Exploitation \textnormal{[SELECT TWO]}}
	
	A bank's EDD team is reviewing the account of the ``Horizon Relief Foundation,'' an NPO registered in a FATF-compliant jurisdiction. The foundation's stated mission is emergency medical aid delivery in conflict zones. The account has received the following transactions over six months: (1) a \$1.2 million wire from a private donor company registered in a FATF Grey List jurisdiction; (2) two cash withdrawals of \$180,000 and \$220,000 at a branch 250 kilometers from the NPO's registered office; (3) transfers of \$600,000 to four sub-NPOs in Conflict Zone Q, where terrorist groups are known to control humanitarian aid distribution channels; (4) an annual audited financial report submitted by the NPO showing \$800,000 in legitimate medical supply purchases with proper invoices and delivery receipts.
	
	Which two transactions or patterns are the most significant red flags requiring SAR consideration?
	
	\begin{enumerate}[label=\Alph*.]
		\item The \$1.2 million wire from a Grey List jurisdiction, combined with the geographic displacement of cash withdrawals (250 km from the registered office) and the lack of documentation linking cash withdrawals to medical supply purchases, constitutes a classic cash diversion pattern associated with terrorist financing exploitation of NPOs.
		\item The transfers to four sub-NPOs in Conflict Zone Q require heightened scrutiny because terrorist groups are known to control aid distribution in the area, which raises the risk that funds designated for medical aid may be diverted to financing terrorist activities rather than reaching beneficiaries.
		\item The audited financial report showing legitimate medical purchases is a red flag because NPOs involved in terrorist financing routinely fabricate audit reports.
		\item The existence of the audited financial report fully mitigates all other red flags identified in the account and confirms that the NPO is operating legitimately.
		\item NPOs operating in conflict zones are inherently exempt from AML monitoring under FATF Recommendation 8 because restricting their banking access would harm humanitarian operations.
	\end{enumerate}
	\textbf{Correct Answers: A and B}
	
	%--- Q40 ---
	\subsection{40. Crowdfunding-Enabled Terrorist Financing \textnormal{[SELECT ONE]}}
	
	An FIU analyst receives a referral from a domestic social media platform. The platform has flagged an account operated by an individual, ``Brother Salim,'' that has posted fundraising appeals describing a ``humanitarian crisis'' in Conflict Zone X, where a UN-designated terrorist organization operates. Over 90 days, 3,400 small donations (\$5 to \$150 each) have been received from 2,800 different donors across 18 countries, totaling approximately \$174,000. No single donation exceeds any standard reporting threshold. The fundraising content does not explicitly mention the terrorist organization but uses imagery and coded language consistent with the group's known propaganda materials. The funds are collected through a payment aggregator and transferred weekly to a cryptocurrency wallet.
	
	What is the primary regulatory concern, and what makes this pattern particularly difficult to detect through standard AML transaction monitoring?
	
	\begin{enumerate}[label=\Alph*.]
		\item The primary concern is tax evasion; the pattern is difficult to detect because the payment aggregator does not file 1099 forms for small donations.
		\item The primary concern is crowdfunding-enabled terrorist financing (TF). Each individual donation falls below reporting thresholds, making the pattern invisible to standard CTR-based detection. The aggregated amount (\$174,000) funds a terrorist-linked purpose effectively, while the decentralized structure of thousands of small donors complicates investigation of the funding network. Standard transaction monitoring is calibrated to flag individual high-value transactions or structuring patterns, not the aggregation of thousands of micro-donations across different senders to a single recipient.
		\item The primary concern is money laundering; the conversion to cryptocurrency at the end constitutes layering.
		\item Because no single transaction exceeds a reporting threshold, no regulatory obligation applies and the pattern need not be investigated further.
		\item The concern is sanctions evasion; the coded language in the appeals constitutes a sanctioned communication under OFAC regulations.
	\end{enumerate}
	\textbf{Correct Answer: B}
	
	%==============================================================
	\section{Part 9: Real Estate, MSBs, Casinos, and Gatekeepers}
	%==============================================================
	
	%--- Q41 ---
	\subsection{41. Real Estate: FinCEN GTO and All-Cash Purchase Red Flags \textnormal{[SELECT ONE]}}
	
	A national title insurance company operating in Miami, Los Angeles, and Manhattan receives a FinCEN Geographic Targeting Order requiring it to identify the natural persons behind legal entities that purchase residential real estate above \$300,000 in all-cash transactions. During the covered period, the title company processes the following transaction: a limited liability company formed in Delaware three weeks prior purchases a \$2.1 million Miami condominium. The wire transfer arrives from a law firm's IOLTA account at a regional bank. The LLC's registered agent is a corporate service provider that declines to disclose the beneficial owner. No mortgage is involved. The seller's closing attorney notes that the buyer's attorney also declined to identify the beneficial owner, citing privilege.
	
	Under the GTO and applicable AML standards, which of the following most accurately describes the title company's obligations?
	
	\begin{enumerate}[label=\Alph*.]
		\item The title company has no AML obligations because it is not a financial institution under the BSA and GTOs are advisory only.
		\item Under the active GTO, the title company is required to identify the natural persons (ultimate beneficial owners) behind the LLC. The use of an IOLTA account, a recently formed anonymous LLC with an uncooperative registered agent, and attorney refusal to identify the beneficial owner all constitute circumstances where the title company cannot satisfy the GTO requirement. The title company should decline to close the transaction until beneficial ownership is disclosed, and must file a FinCEN currency transaction report (or equivalent GTO report form) identifying the suspicious circumstances if the transaction proceeds.
		\item The title company's obligation is satisfied by obtaining the LLC's certificate of formation from Delaware, which is publicly available.
		\item The attorney-client privilege invoked by the buyer's attorney overrides the GTO's beneficial ownership identification requirement for real estate attorneys.
		\item The GTO only applies to all-cash transactions above \$5 million; the \$2.1 million transaction is below this threshold.
	\end{enumerate}
	\textbf{Correct Answer: B}
	
	%--- Q42 ---
	\subsection{42. Casino: Chip Wash and SAR Obligations \textnormal{[SELECT ONE]}}
	
	Over a single weekend at a Las Vegas casino, a player identified as Mr. G purchases \$300,000 in chips using cash delivered in a duffel bag. Mr. G plays blackjack for approximately 45 minutes with \$1,500 in chips (the remainder sitting on the table), then cashes out \$297,000 in chips for a casino check payable to his name. He accepts the \$3,000 ``gaming loss'' as the cost of the transaction. The casino's AML compliance officer is reviewing whether to file a SAR and a CTR.
	
	Which of the following most accurately describes the casino's obligations?
	
	\begin{enumerate}[label=\Alph*.]
		\item The casino's only obligation is to file a CTR for the \$300,000 cash purchase of chips. No SAR is required because Mr. G engaged in actual gaming.
		\item The casino must file a CTR for the \$300,000 cash-in transaction, because it exceeds the \$10,000 gaming-day threshold. The casino must also file a SAR because the transaction pattern -- purchasing a large volume of chips with cash, engaging in minimal gaming, and cashing out for a check -- constitutes a classic ``chip washing'' scheme where the casino's financial system is used to convert criminal cash into a negotiable casino check, effectively laundering the funds through the legitimizing appearance of gaming activity.
		\item Because Mr. G actually gambled and lost \$3,000, the transaction has a legitimate gaming purpose and neither a CTR nor a SAR is required.
		\item The casino should file a SAR only if Mr. G's true identity cannot be verified; a CTR is not required for chip purchases, only for cash-out transactions.
		\item Chip wash transactions at casinos are not covered by the BSA because gaming chips are not legal tender.
	\end{enumerate}
	\textbf{Correct Answer: B}
	
	%--- Q43 ---
	\subsection{43. Hawala Network: Identifying the Transaction Cycle and Regulatory Obligation \textnormal{[SELECT TWO]}}
	
	A financial crimes detective is investigating a hawala network operating across five U.S. cities. Undercover agents document the following: hawaladar A in New York accepts \$80,000 in cash from a drug trafficker's courier, records a coded reference number, and calls hawaladar B in Karachi, Pakistan. Hawaladar B pays the equivalent in Pakistani Rupees (at a below-market exchange rate favorable to the hawaladar) to a designated recipient within 24 hours. No wire transfer occurs; no bank account is debited or credited; no financial institution processes the transaction. Hawaladar A operates from a convenience store and has no FinCEN MSB registration. Over 36 months, the network processed an estimated \$22 million.
	
	Which two of the following are accurate regulatory and investigative findings?
	
	\begin{enumerate}[label=\Alph*.]
		\item Hawaladar A is required to register as a Money Services Business with FinCEN under FATF Recommendation 14 and the BSA, because the operation of an informal value transfer system qualifies as money transmitting regardless of whether a bank account is used.
		\item Hawaladar A's failure to register as an MSB, failure to implement AML controls, and processing of criminal proceeds may expose him to criminal liability under 18 U.S.C.\ \S 1960 (operating an unlicensed money transmitting business) and separate money laundering charges under 18 U.S.C.\ \S 1956.
		\item Hawaladar A has no registration or AML obligation because the transaction never touches the U.S. banking system and therefore falls outside BSA jurisdiction.
		\item The hawala transaction constitutes placement and layering simultaneously, because criminal cash is being removed from the U.S. and converted into a foreign currency obligation with no documentary trail.
		\item The hawala system is exempt from AML regulation under all FATF member countries because of its cultural heritage status as a traditional value transfer mechanism.
	\end{enumerate}
	\textbf{Correct Answers: A and B}
	
	%--- Q44 ---
	\subsection{44. Gatekeeper Risk: Law Firm IOLTA Account as Layering Vehicle \textnormal{[SELECT ONE]}}
	
	A law firm receives \$5.5 million in a single wire transfer into its IOLTA pooled trust account, purportedly representing the settlement proceeds of a commercial litigation matter on behalf of a foreign corporate client. The wire originates from a shell company in Cyprus. The law firm's own retainer agreement with the client covers only a \$12,000 advisory engagement. The partner handling the matter has never met the client in person; all instructions have been received via encrypted messaging. The law firm withdraws the full \$5.5 million within five business days, transfers \$5.2 million to a second shell company in the UAE (per client instruction), and retains \$300,000 as a ``success fee.'' The originating bank flags the wire and contacts the law firm's bank.
	
	What makes the law firm's IOLTA account particularly dangerous as a money laundering instrument in this scenario?
	
	\begin{enumerate}[label=\Alph*.]
		\item IOLTA accounts are subject to the highest level of BSA scrutiny because law firms are required to file CTRs on all pooled account transactions exceeding \$10,000.
		\item The law firm's IOLTA account functions as a blind pass-through vehicle. The receiving bank typically extends professional trust to the law firm and does not look through the pooled account to examine individual client fund flows. This allows illicit funds to enter the banking system credentialed by an attorney's professional status, commingled with legitimate client funds, and disbursed quickly to offshore shell companies without triggering standard transaction monitoring controls. The disproportionate amount relative to the retainer, the offshore shell company source, the absence of face-to-face contact, and the rapid disbursement are all red flags that the law firm may be -- knowingly or unknowingly -- facilitating money laundering.
		\item IOLTA accounts are exempt from all AML reporting obligations under attorney-client privilege, which overrides the BSA.
		\item The danger lies in the large success fee, which suggests the law firm is charging above-market rates.
		\item IOLTA accounts cannot be the subject of a SAR filing because only bank accounts held directly by customers can generate SAR obligations.
	\end{enumerate}
	\textbf{Correct Answer: B}
	
	%==============================================================
	\section{Part 10: Advanced Scenario Analysis}
	%==============================================================
	
	%--- Q45 ---
	\subsection{45. Scenario: The Multi-Layered Sovereign Corruption Extraction \textnormal{[SELECT ONE]}}
	
	\textbf{Background:} The Prime Minister of Country Zeta, a resource-rich nation with a Transparency International Corruption Perception Index score of 12/100, has been diverting state oil revenues for 11 years. The scheme works as follows: State-owned oil company OilZeta signs ``service agreements'' with a British Virgin Islands shell company (BVI-Co) for ``logistics consulting.'' BVI-Co receives \$220 million in wire transfers directly from OilZeta's account at a Swiss private bank. BVI-Co's funds are then split: \$150 million is wired to a Cayman Islands trust of which the Prime Minister's adult son is the sole beneficiary; \$70 million is invested in London commercial real estate through a UK LLP whose members are two more BVI shell companies. The Prime Minister has never held a personal account at any of the institutions involved.
	
	\textbf{Question:} A new compliance officer at the Swiss private bank has inherited the OilZeta account. During a routine file review, she identifies the following: OilZeta is a state-owned enterprise; no counterparty analysis was performed on BVI-Co at account opening; the service agreements lack any technical deliverables or performance milestones; and the Prime Minister's son has been publicly named in a foreign court filing as a beneficiary of a trust receiving funds from state oil revenues. What is her most critical immediate obligation?
	
	\begin{enumerate}[label=\Alph*.]
		\item To close the OilZeta account immediately and return all funds to OilZeta's home jurisdiction, because continuing the relationship creates regulatory liability.
		\item To escalate the findings to senior management and the compliance committee immediately, classifying OilZeta and its associated payment flows as a high-risk PEP-adjacent relationship (given the state ownership connection to the Prime Minister). She must initiate an emergency EDD review covering the source of OilZeta's funds, the nature of the BVI-Co service agreements, and the beneficial ownership of all downstream entities receiving wires. Given the court filing linking the Prime Minister's son to the trust, she must also assess whether the transactions constitute proceeds of corruption, and file STRs with the Swiss FIU if the facts support a reasonable suspicion of money laundering. She cannot close the account without senior management approval and must ensure no action tips off the account holder.
		\item To wait for the foreign court proceeding to conclude before taking any action, because STR filings are premature until a conviction is secured.
		\item To contact the Prime Minister of Country Zeta directly to obtain a source of funds explanation for the BVI-Co service fees.
		\item To request that BVI-Co provide updated service agreements with performance milestones and, once received, consider the matter resolved.
	\end{enumerate}
	\textbf{Correct Answer: B}
	
	%--- Q46 ---
	\subsection{46. Scenario: The Compromised Correspondent Bank Trap (FATF Effectiveness Context) \textnormal{[SELECT TWO]}}
	
	\textbf{Background:} DeltaBank is a correspondent bank in the EU providing USD clearing to SouthernTrade Bank (STB), a commercial bank in Country Y. Country Y received a FATF rating of ``Non-Compliant'' on Technical Compliance Recommendation 10 (CDD) and ``Low'' on Effectiveness IO.4 (Supervisory oversight). STB's CBDDQ discloses that it services 8 unlicensed money changers, 3 casino-affiliated entities, and offers payable-through access to 11 undisclosed sub-correspondents. DeltaBank's relationship with STB generates   800,000 in annual fee income. Over 18 months, DeltaBank's transaction monitoring generates 34 alerts on STB's account, all of which are closed by a single analyst with the note ``STB is a licensed bank.'' No STRs have been filed. An external audit rates DeltaBank's monitoring of the STB account as ``inadequate.''
	
	Which two of the following represent the most serious findings that DeltaBank's compliance committee must address?
	
	\begin{enumerate}[label=\Alph*.]
		\item The 34 alerts closed without adequate investigation constitute a systemic monitoring failure. A licensed bank status is not a sufficient reason to close an alert; each alert required individual analysis of the underlying transaction patterns, and the analyst's blanket closure methodology is a material AML control deficiency.
		\item STB's CBDDQ profile -- servicing unlicensed money changers, casino entities, and 11 undisclosed sub-correspondents in a Non-Compliant jurisdiction -- means DeltaBank's risk assessment of the relationship was either not completed or is materially inadequate. This profile should have triggered EDD, potential business restrictions, and documented senior management review of whether to maintain the relationship.
		\item The euro 800,000 in annual fee income means the relationship is commercially significant and cannot be terminated without board approval; this justifies the current monitoring approach.
		\item The external audit's finding of ``inadequate'' monitoring is an opinion and cannot be used as a basis for regulatory action against DeltaBank.
		\item DeltaBank's liability is limited because the monitoring alerts were reviewed (even if inadequately) and the analyst documented his rationale.
	\end{enumerate}
	\textbf{Correct Answers: A and B}
	
	%--- Q47 ---
	\subsection{47. Scenario: The Money Mule Network at a Fintech P2P Platform \textnormal{[SELECT ONE]}}
	
	\textbf{Background:} A fintech company operates a peer-to-peer lending platform. Its AML team identifies the following cluster of accounts: 17 accounts all registered within a 48-hour window, all using IP addresses that route through a VPN service associated with a known secrecy jurisdiction. All 17 accounts received small P2P loans (\$2,000--\$5,000 each) funded by a single institutional lender account on the platform. Each of the 17 accounts repaid the loans within 3 days using wire transfers from offshore accounts at three different island jurisdiction banks. After repayment, all 17 accounts initiated withdrawals to separate prepaid debit card numbers. The institutional lender account belongs to ``Digital Micro-Finance Solutions LLC,'' registered in a state with no beneficial ownership disclosure requirements, and it has funded over 400 similar micro-loans in the past six months.
	
	What is the most likely criminal typology being executed through this platform, and what are the key AML indicators?
	
	\begin{enumerate}[label=\Alph*.]
		\item This is a standard peer-to-peer lending fraud scheme. The key indicator is the rapid repayment, which is consistent with short-term arbitrage strategies common in fintech platforms.
		\item This is a self-funding layering ring exploiting the P2P lending infrastructure as a laundering mechanism. The institutional lender (likely controlled by the criminal network) issues loans to mule accounts it also controls; the ``loan repayment'' from offshore accounts converts offshore illicit funds into platform-processed domestic transactions that appear as legitimate loan repayments; the withdrawal to prepaid cards completes the integration stage. Key AML indicators: clustered account registration, common VPN routing, offshore loan repayment sources, prepaid card withdrawals, and a high-volume anonymous institutional lender with no transparent ownership.
		\item The scheme is a form of trade-based money laundering because the platform processes commercial transactions between the institutional lender and the borrowers.
		\item The scheme is most accurately characterized as structuring, because each loan is below the CTR threshold.
		\item The fintech platform has no AML obligation because it is a technology platform rather than a bank, and loan-based transactions are not covered by BSA.
	\end{enumerate}
	\textbf{Correct Answer: B}
	
	%--- Q48 ---
	\subsection{48. Scenario: Human Trafficking -- Financial Profiling of a Controller Account \textnormal{[SELECT TWO]}}
	
	\textbf{Background:} A regional bank's transaction monitoring flags an account held by Ms. Wang, listed as the owner of a licensed nail salon. The account receives between \$18,000 and \$24,000 per month in cash deposits, all made in amounts between \$700 and \$1,400. Additionally, the account receives weekly Venmo and Zelle payments from 28 different individuals. Every Monday, approximately 90\% of the balance is withdrawn in cash. The account also pays utility bills, mobile phone plans, and internet service for six different residential addresses. Public records indicate that all six addresses are zoned for multi-family residential use and are registered to two LLCs whose registered agent is the same person as Ms. Wang's nail salon's registered agent. The nail salon's business profile does not explain the volume or pattern of transactions.
	
	Which two of the following represent the most significant indicators of potential human trafficking in this account profile?
	
	\begin{enumerate}[label=\Alph*.]
		\item The pattern of receiving payments from 28 different individuals and paying utility expenses for six residential addresses in the same controller's network is consistent with a labor or sex trafficking operation, where workers' earnings are deposited to a single controller account and living expenses for the controlled housing are paid by the controller to maintain physical and financial control over the victims.
		\item The Monday cash withdrawal pattern (90\% of weekly accumulation) suggests the controller is physically collecting and redistributing funds on a weekly cycle, consistent with paying ``wages'' to traffickers or transmitting proceeds to a higher-level criminal network.
		\item The nail salon business fully explains the account activity because beauty service businesses legitimately receive cash deposits and P2P payments from individual customers.
		\item The LLC ownership of the residential properties is a standard small business real estate investment strategy and is not an AML indicator.
		\item The six utility payments indicate the account holder is generously paying utility bills for family members, which is a common cultural practice with no AML significance.
	\end{enumerate}
	\textbf{Correct Answers: A and B}
	
	%--- Q49 ---
	\subsection{49. Scenario: De-Risking Decision -- Individual Assessment vs.\ Blanket Termination \textnormal{[SELECT ONE]}}
	
	\textbf{Background:} A global tier-1 bank has received a compliance mandate from its group risk committee to ``terminate all correspondent banking relationships with banks headquartered in Sub-Saharan African countries'' due to the region's ``elevated AML risk.'' The mandate was issued following a regulatory action at a competitor institution involving a single African correspondent bank. The group's compliance department has not conducted individual risk assessments of its 22 African correspondent relationships. Several of these 22 banks have been individually assessed by their local regulators as having robust AML programs, including three that have passed recent FATF mutual evaluations with ``Compliant'' or ``Largely Compliant'' ratings.
	
	What is the most significant AML and regulatory risk created by this blanket termination mandate?
	
	\begin{enumerate}[label=\Alph*.]
		\item The mandate is legally permissible because banks have absolute discretion to choose their correspondent partners and can terminate any relationship without providing justification.
		\item The blanket termination mandate violates the risk-based approach to AML, which requires institutions to conduct individual assessments rather than apply categorical exclusions. By terminating compliant banks along with non-compliant ones, the mandate: (i) deprives legitimate financial institutions in 22 countries of access to USD clearing, causing financial exclusion harm; (ii) may drive their customers to unregulated, informal financial channels that are less transparent; (iii) may expose the group to criticism from FATF, financial supervisors, and relevant government authorities for implementing a policy that contradicts the RBA principles; and (iv) may create legal liability under anti-discrimination frameworks in some jurisdictions.
		\item The mandate creates no AML risk; banks that terminate high-risk relationships always reduce their ML exposure.
		\item The mandate is appropriate because the risk-based approach allows banks to apply regional-level risk assessments without individual entity analysis.
		\item The primary risk is reputational; there is no regulatory or legal concern with a bank choosing its business partners.
	\end{enumerate}
	\textbf{Correct Answer: B}
	
	%--- Q50 ---
	\subsection{50. Scenario: Post-SAR Filing Management and Regulatory Exam Readiness \textnormal{[SELECT TWO]}}
	
	\textbf{Background:} A bank's AML team has filed 7 SARs over 14 months on the accounts of three related corporate entities: OceanTrade Ltd., MarineShip LLC, and PortLogis Inc. The SARs document a pattern of structured cash deposits, rapid wire transfers to offshore accounts in high-risk jurisdictions, and invoice discrepancies in associated trade finance transactions consistent with TBML. The bank continues to maintain all three accounts at the request of the lead account relationship manager, who argues that the customers are ``commercially important.'' No formal decision document exists for the continued account maintenance. The bank's BSA Officer has not been consulted on the accounts since the second SAR was filed 10 months ago.
	
	Which two of the following represent the most critical compliance failures in the bank's management of this situation?
	
	\begin{enumerate}[label=\Alph*.]
		\item Continuing to maintain accounts after multiple SAR filings without a documented, BSA Officer-approved decision supported by a risk rationale is a material compliance failure. Ongoing SAR-subject accounts require formal senior management review, documented justification for continued service, and heightened ongoing monitoring.
		\item The BSA Officer's exclusion from account management decisions after the second SAR filing is a governance failure. The BSA Officer must be involved in decisions about continuing relationships with accounts on which multiple SARs have been filed, because maintaining such accounts without BSA Officer knowledge and approval undermines the independence and authority of the AML compliance function.
		\item The relationship manager's commercial importance argument is a valid AML risk management consideration that can override the standard SAR-follow-up review process.
		\item Seven SARs on the same entities over 14 months is evidence that the transaction monitoring system is working correctly; no additional steps are required.
		\item The bank's only obligation after filing a SAR is to file additional SARs if new suspicious activity is identified; no formal account management review process is required by law.
	\end{enumerate}
	\textbf{Correct Answers: A and B}
	
	%==============================================================
	\section{Part 11: Regulatory Frameworks, Governance, and Enforcement}
	%==============================================================
	
	%--- Q51 ---
	\subsection{51. Independent Testing: Scope and Staffing Requirements \textnormal{[SELECT ONE]}}
	
	An investment bank's audit committee is reviewing a proposal to have its BSA/AML independent testing conducted by the bank's internal audit department. The internal audit department reports to the Chief Audit Executive, who reports to the Audit Committee of the Board of Directors. However, the internal audit team that would conduct the AML test is the same team that: (1) helped select the AML software vendor two years ago; (2) co-authored the current AML policy manual alongside the compliance team; and (3) has three members who previously worked as AML analysts in the compliance department. The Compliance Officer argues that because internal audit reports to the board, it qualifies as fully independent.
	
	What is the critical flaw in the proposed arrangement, and what must be done to satisfy the independence requirement?
	
	\begin{enumerate}[label=\Alph*.]
		\item The arrangement fully satisfies independence requirements because the reporting line to the board is the only independence criterion under the BSA.
		\item The arrangement is flawed because independence is not solely a matter of organizational reporting lines. Members of the testing team who participated in designing, building, or staffing the controls being tested cannot independently audit those controls. To satisfy independence, the bank must either use an external auditor or an internal audit team whose members have no prior involvement in designing, selecting, or operating the AML controls being tested. The prior co-authorship of the AML policy, vendor selection participation, and former compliance department roles of the proposed auditors create material independence impairments.
		\item The arrangement is flawed only with respect to the three former AML analysts; the remaining audit team members can conduct the test independently.
		\item The arrangement satisfies independence as long as the prior involvement of audit team members occurred more than two years ago.
		\item Internal audit can never conduct AML independent testing under any circumstances; only a licensed external auditor is acceptable under BSA regulations.
	\end{enumerate}
	\textbf{Correct Answer: B}
	
	%--- Q52 ---
	\subsection{52. Whistleblower Protections and Retaliation \textnormal{[SELECT ONE]}}
	
	A compliance analyst at a regional bank submits an internal report to the bank's ethics hotline documenting what she believes is a systematic failure by senior management to investigate SAR-eligible transactions involving a high-revenue corporate client. Three weeks after submitting the report, the analyst receives a performance improvement plan (PIP) citing ``communication issues'' -- her first negative performance review in six years of employment. She believes the PIP is retaliatory. Under U.S. law, which regulatory framework and agency most directly addresses her situation?
	
	\begin{enumerate}[label=\Alph*.]
		\item The bank's internal HR policy is the only recourse available; there is no federal whistleblower protection for AML compliance analysts.
		\item The analyst is protected under the anti-retaliation provisions of the Financial Institutions Reform Recovery and Enforcement Act (FIRREA), the Bank Secrecy Act (31 U.S.C.\ \S 5328), and potentially the Dodd-Frank Act Section 922 (if the disclosure is SEC-related). The BSA specifically prohibits financial institutions from retaliating against employees who report violations of the BSA. The analyst may file a complaint with the U.S. Department of Labor (OSHA), which investigates BSA whistleblower retaliation claims, within 180 days of the adverse employment action.
		\item The analyst is protected only if she filed the report externally with FinCEN; internal hotline reports are not covered by federal whistleblower protections.
		\item The Sarbanes-Oxley Act (SOX) is the primary protection, and the analyst must file her complaint within 30 days of the retaliation.
		\item There is no federal agency with jurisdiction over AML whistleblower retaliation at a non-publicly listed bank.
	\end{enumerate}
	\textbf{Correct Answer: B}
	
	%--- Q53 ---
	\subsection{53. Recordkeeping: Five-Year Rule and Documentary Requirements \textnormal{[SELECT TWO]}}
	
	A compliance director is conducting a documentation retention review in advance of a bank regulatory examination. She identifies five categories of records and needs to confirm which retention requirements apply. The five categories are: (1) CTR filings; (2) SAR filings and supporting documentation; (3) Customer identification program documents (CIP); (4) Correspondent bank certification records under Section 319(b); (5) Wire transfer records for transactions of \$3,000 or more.
	
	Which two of the following retention rules are correct?
	
	\begin{enumerate}[label=\Alph*.]
		\item SAR filings and all supporting documentation (Category 2) must be retained for five years from the date of filing under 31 U.S.C.\ \S 5318(g).
		\item Wire transfer records for transactions of \$3,000 or more (Category 5) must be retained for five years from the date of the transaction under 31 CFR \S 1020.410.
		\item CTR filings (Category 1) must be retained for only one year from the date of filing; only SARs require the five-year retention period.
		\item CIP documents (Category 3) may be discarded once the customer relationship has been terminated for more than 12 months.
		\item Correspondent bank certification records (Category 4) need not be retained because they are the respondent bank's document rather than the correspondent bank's own record.
	\end{enumerate}
	\textbf{Correct Answers: A and B}
	
	%--- Q54 ---
	\subsection{54. GDPR and AML: Navigating Conflicting Legal Obligations \textnormal{[SELECT ONE]}}
	
	A French bank's data protection officer (DPO) and its BSA/AML compliance officer are in disagreement. A customer submits a GDPR ``right to erasure'' (right to be forgotten) request, demanding the bank delete all personal data held about them, including KYC documents, transaction records, and a SAR filed on the customer. The DPO argues that the GDPR right to erasure is absolute and the bank must delete all personal data within 30 days. The compliance officer argues that AML/CFT recordkeeping obligations override the GDPR request.
	
	Which of the following most accurately resolves the conflict?
	
	\begin{enumerate}[label=\Alph*.]
		\item The DPO is correct; GDPR is EU primary law and takes absolute precedence over member state AML legislation.
		\item The compliance officer is correct that GDPR Article 17(3) explicitly carves out exceptions to the right to erasure where retention is necessary for compliance with a legal obligation under EU or member state law. AML/CFT recordkeeping obligations (under the EU Anti-Money Laundering Directives and national implementing legislation) constitute such a legal obligation. The bank may lawfully refuse to erase records it is legally required to retain for AML/CFT purposes, including KYC, transaction records, and SAR documentation. Furthermore, the SAR confidentiality obligation independently prohibits acknowledging the existence of the SAR to the customer.
		\item The bank must delete all records except the SAR itself; the SAR is the only document exempt from GDPR erasure.
		\item Neither obligation takes precedence; the bank must escalate the conflict to the European Data Protection Board for a formal determination.
		\item The right to erasure applies to all data held for more than five years; any data within the five-year AML retention window is exempt, but data beyond five years must be deleted.
	\end{enumerate}
	\textbf{Correct Answer: B}
	
	%--- Q55 ---
	\subsection{55. Basel Committee: Three Lines of Defense Application to AML \textnormal{[SELECT TWO]}}
	
	A large international bank is implementing the Basel Committee on Banking Supervision's guidance on the ``Sound Management of Risks Related to ML and TF'' (2014 and revised 2017). The bank's management is debating how to structure accountability for AML risk management across its three lines of defense. The following proposals are made: (i) business line managers (First Line) should independently determine AML risk assessments for their own products without compliance oversight; (ii) the AML compliance function (Second Line) should have authority to approve or reject customer relationships based on risk assessments; (iii) Internal Audit (Third Line) should periodically test the effectiveness of First and Second Line AML controls and report findings to the Audit Committee, independent of management.
	
	Which two of the following assessments of the proposals are correct?
	
	\begin{enumerate}[label=\Alph*.]
		\item Proposal (i) is flawed. Business line managers (First Line) own and manage risks but must operate within an AML risk management framework set by the Second Line (compliance). Allowing the First Line to independently determine risk assessments without compliance oversight undermines the risk framework design function of the Second Line.
		\item Proposal (iii) is consistent with Basel Committee guidance. Internal audit, as the Third Line, must independently assess whether AML controls throughout the institution are effective and must report its findings directly to the Audit Committee or Board, free from management interference.
		\item Proposal (ii) is flawed because the compliance function should never be involved in customer relationship decisions; all onboarding decisions are exclusively a First Line responsibility.
		\item Proposal (i) is consistent with Basel guidance because business line managers have the best knowledge of their own product risks.
		\item Proposal (iii) is flawed because Internal Audit should report to the Chief Compliance Officer, not to the Audit Committee, to ensure effective coordination.
	\end{enumerate}
	\textbf{Correct Answers: A and B}
	
	%==============================================================
	\section{Part 12: Final High-Difficulty Mixed Questions}
	%==============================================================
	
	%--- Q56 ---
	\subsection{56. Correspondent Banking: Reverse Correspondent Risk and Shell Bank Detection \textnormal{[SELECT ONE]}}
	
	A small bank in Country Q (a developing economy) maintains a USD correspondent account at a large U.S. clearing bank, FirstClearing. Recently, FirstClearing discovers that the Country Q bank has been processing USD transactions on behalf of an entity that FirstClearing has never heard of: ``Global Settlement Nominees Inc.,'' which appears to operate exclusively through the Country Q bank's account. Investigation reveals that Global Settlement Nominees has no verifiable physical presence, no licensed banking activities in any jurisdiction, no employees, and processes an average of \$90 million per month in USD transactions, far exceeding the Country Q bank's own expected transaction volume. The entity cannot be found in any public corporate registry.
	
	What is the most accurate characterization of Global Settlement Nominees under the USA PATRIOT Act, and what must FirstClearing do?
	
	\begin{enumerate}[label=\Alph*.]
		\item Global Settlement Nominees appears to be a foreign shell bank as defined under Section 313 of the USA PATRIOT Act (a bank with no physical presence in any jurisdiction and not affiliated with a regulated depository institution). FirstClearing is absolutely prohibited from maintaining correspondent account access for a foreign shell bank. It must immediately terminate the arrangement that provides Global Settlement Nominees with access to its payment infrastructure through the Country Q bank, file a SAR documenting the finding, and review whether the Country Q bank itself should continue as a respondent.
		\item Global Settlement Nominees is simply an undisclosed corporate banking client of the Country Q bank. Because it is the Country Q bank's customer (not FirstClearing's), FirstClearing has no independent obligation to investigate or act.
		\item FirstClearing should file a SAR but continue the relationship to avoid disrupting legitimate financial services in Country Q.
		\item FirstClearing's obligations are met by increasing its monitoring frequency for the Country Q bank's account from annual to semi-annual review.
		\item Global Settlement Nominees is exempt from shell bank rules because it is technically not a bank; shell bank prohibitions only apply to licensed banking entities, not to corporate payment processors.
	\end{enumerate}
	\textbf{Correct Answer: A}
	
	%--- Q57 ---
	\subsection{57. Cross-Border Currency Movement: Bulk Cash Smuggling vs.\ Structuring \textnormal{[SELECT ONE]}}
	
	A customs intelligence unit at a major international airport has identified two concurrent patterns during baggage screening operations:
	
	\textbf{Pattern A:} A single traveler is found to be carrying \$185,000 in cash concealed in a specially constructed false-bottom suitcase. The traveler did not declare the currency on the customs declaration form.
	
	\textbf{Pattern B:} Over the same 60-day period, 22 different travelers on the same airline route each carried between \$8,500 and \$9,800 in cash and individually declared their currency. Each declaration is technically accurate for the amounts individually carried, but no single traveler declared over \$10,000. All 22 travelers share the same home address (a mail drop) and all conducted transactions at the same bank branch in the 48 hours prior to departure.
	
	Which of the following most accurately distinguishes the AML violations in each pattern?
	
	\begin{enumerate}[label=\Alph*.]
		\item Only Pattern A is a violation; Pattern B involves accurate declarations and does not constitute any illegal activity.
		\item Pattern A constitutes bulk cash smuggling under 31 U.S.C.\ \S 5332 (failure to declare currency exceeding \$10,000 at the border). Pattern B constitutes smurfing / structuring, where a network of individuals has coordinated to fragment a large currency movement across 22 carriers, each carrying sub-threshold amounts. While each declaration is individually accurate, the aggregate movement and coordination evidence a deliberate scheme to circumvent border currency reporting requirements. Both are federal crimes; Pattern B may also implicate money mule recruitment and conspiracy charges.
		\item Both patterns are violations of the same statute (31 U.S.C.\ \S 5316); the only distinction is the method of concealment.
		\item Pattern B is not a violation because all declarations were accurate; the shared home address and common bank branch are coincidences that do not establish criminal intent.
		\item Pattern A constitutes tax evasion; Pattern B constitutes customs fraud. Neither is a money laundering violation under the BSA.
	\end{enumerate}
	\textbf{Correct Answer: B}
	
	%--- Q58 ---
	\subsection{58. Predicate Offense: Self-Laundering and Double Jeopardy \textnormal{[SELECT TWO]}}
	
	A federal prosecutor is considering charging Mr. R with both wire fraud (the predicate offense) and money laundering under 18 U.S.C.\ \S 1956 for the subsequent movement of the fraud proceeds through a shell company network. Mr. R's defense attorney argues two points: (1) that charging both the predicate offense (wire fraud) and the money laundering of the fraud proceeds constitutes double jeopardy because the same underlying conduct is being charged twice; and (2) that self-laundering (where the predicate offender also launders their own proceeds) is not covered under U.S. money laundering statutes.
	
	Which two of the following most accurately assess the defense attorney's arguments?
	
	\begin{enumerate}[label=\Alph*.]
		\item The double jeopardy argument is without merit. A money laundering charge and a predicate offense charge involve legally distinct elements and distinct criminal acts (committing the underlying crime vs.\ conducting financial transactions with the proceeds). Prosecuting both simultaneously does not violate the Double Jeopardy Clause because each offense requires proof of a different element.
		\item The self-laundering argument is without merit under U.S. law. 18 U.S.C.\ \S 1956 explicitly applies to persons who conduct financial transactions with the proceeds of specified unlawful activity with intent to conceal, promote, or evade taxes -- it does not require a separate third-party launderer. The predicate offender can also be charged as the money launderer.
		\item The double jeopardy argument is correct; charging both wire fraud and money laundering for the same conduct violates the Fifth Amendment.
		\item The self-laundering argument is correct; FATF Recommendation 3 permits but does not require countries to criminalize self-laundering, and the U.S. has opted out of applying this standard.
		\item Both of the defense attorney's arguments are legally sound and the prosecutor must choose one charge or the other.
	\end{enumerate}
	\textbf{Correct Answers: A and B}
	
	%--- Q59 ---
	\subsection{59. AI-Based Transaction Monitoring: Model Risk and Explainability \textnormal{[SELECT ONE]}}
	
	A bank has deployed a machine learning-based transaction monitoring system that uses a deep neural network model. The model demonstrates an 82\% accuracy rate in identifying true positive suspicious activity, compared to 34\% for the prior rules-based system. However, the model operates as a ``black box'' -- compliance analysts are unable to understand why specific transactions are flagged or why specific accounts are rated as high-risk. The model's developer cannot explain the decision-making logic in plain language. A banking regulator conducts an examination and rates the model governance as ``unsatisfactory.''
	
	Why is the model's high accuracy rate insufficient to satisfy the regulator's governance concerns?
	
	\begin{enumerate}[label=\Alph*.]
		\item The regulator is incorrect; accuracy rate is the only metric that matters for AML transaction monitoring systems, and an 82\% accuracy rate should be rated satisfactory.
		\item Model explainability is a core governance requirement for AML transaction monitoring systems. Compliance analysts must be able to understand why a transaction was flagged in order to complete a meaningful investigation and determine whether to file a SAR. An opaque model creates several risks: (i) analysts cannot identify and correct model biases or errors; (ii) the model cannot be validated for effectiveness over time; (iii) regulators cannot assess whether the model is producing alerts on legally sound grounds; and (iv) SAR filings based on unexplained model outputs may not withstand legal scrutiny. High accuracy does not substitute for the governance, validation, and explainability requirements that apply to all AML control systems.
		\item The bank should simply file a SAR for every transaction the model flags, regardless of the analyst's ability to understand the rationale, to satisfy the accuracy-based standard.
		\item The regulator's concern can be resolved by increasing the model's accuracy rate to 95\% through additional training data.
		\item The explainability requirement applies only to sanctions screening models; transaction monitoring models are not subject to the same governance standard.
	\end{enumerate}
	\textbf{Correct Answer: B}
	
	%--- Q60 ---
	\subsection{60. Adverse Media and OSINT: Escalation Triggers \textnormal{[SELECT ONE]}}
	
	A bank's enhanced due diligence team is conducting adverse media screening on a corporate customer, Meridian Exports Ltd. The following findings are discovered through credible investigative journalism databases, government procurement records, and corporate registry databases: (1) Meridian's CEO was named (not charged, not convicted) in a foreign parliamentary inquiry into a government procurement scandal 18 months ago; (2) a company with a similar name and the same registered address was previously dissolved following a regulatory investigation into trade fraud (different jurisdiction); (3) Meridian's primary business counterparty is a company whose sole director appears on the Interpol Red Notice database. Meridian's account has been active for four years with no prior alerts.
	
	What is the correct AML response to these OSINT findings?
	
	\begin{enumerate}[label=\Alph*.]
		\item Because no conviction has occurred and the account has been clean for four years, the adverse media is inconclusive and can be noted in the file without further action.
		\item The three OSINT findings collectively constitute credible adverse media requiring immediate escalation to the compliance committee. The CEO's naming in a parliamentary corruption inquiry (regardless of whether charges were filed), the dissolved predecessor company's trade fraud history, and the Interpol-flagged business counterparty are material risk indicators that may require: re-assessment of Meridian's risk rating (likely elevation to high risk), initiation of an EDD review, enhanced ongoing monitoring, and potentially a SAR filing if the transaction profile supports reasonable suspicion. OSINT findings do not require a conviction to be actionable -- reasonable suspicion of ML/TF risk, not proof of guilt, is the standard.
		\item The Interpol Red Notice on the counterparty's director is sufficient alone to require immediate account closure; the other findings are secondary.
		\item The bank should contact Meridian's CEO directly to obtain an explanation for the parliamentary inquiry finding before taking any compliance action.
		\item The dissolved predecessor company finding is irrelevant because it was a separate legal entity; corporate ancestry has no AML significance.
	\end{enumerate}
	\textbf{Correct Answer: B}
	
	%--- Q61 ---
	\subsection{61. Customer Risk Profile: Integrating Multiple Risk Dimensions \textnormal{[SELECT TWO]}}
	
	A bank is assessing the aggregate risk score for five prospective customers using a multi-dimensional risk matrix with four risk factors: (1) Customer type/PEP status, (2) Geographic risk, (3) Product/service risk, and (4) Source of funds complexity. Each factor is rated Low (1), Medium (2), or High (3). The aggregate score triggers CDD treatment: scores 4--6: Standard CDD; 7--9: Medium EDD; 10--12: Full EDD + Senior Management Approval.
	
	Customer profiles:
	\begin{itemize}
		\item Customer A: Non-PEP domestic individual; domestic jurisdiction (FATF-compliant); savings account; salary deposit from employer. Scores: 1+1+1+1 = 4.
		\item Customer B: Non-PEP foreign national; country on FATF Grey List; wire transfer account; source of funds is ``trading income'' with no documentation. Scores: 1+3+2+3 = 9.
		\item Customer C: PEP (foreign senior government official); country with Corruption Index below 20/100; private banking / trust account; funds from a state enterprise ``consulting fee.'' Scores: 3+3+3+3 = 12.
		\item Customer D: Non-PEP domestic charity; domestic jurisdiction; basic savings account; charitable donation income. Scores: 1+1+1+1 = 4.
		\item Customer E: Former PEP (retired from public office 3 years ago); jurisdiction under FATF Black List; offshore corporate account; source stated as real estate sale proceeds but documentation not provided. Scores: 2+3+3+3 = 11.
	\end{itemize}
	
	Which two of the following statements are correct?
	
	\begin{enumerate}[label=\Alph*.]
		\item Customer C (score 12) and Customer E (score 11) both require Full EDD and Senior Management Approval before account opening, and Customer C's PEP status mandates that the source of consulting fees be independently verified and documented, given the high corruption risk of the home jurisdiction.
		\item Customer B (score 9) falls in the Medium EDD tier. The undocumented ``trading income'' source of funds must be resolved before proceeding; the Grey List jurisdiction elevates risk further, and without documentation of source of funds the account should not be opened.
		\item Customer D (score 4) should be classified as Standard CDD, but because it is a charity with potential NPO-related terrorist financing risk, it should automatically receive Full EDD.
		\item Customer A (score 4) should receive EDD because domestic individuals who use salary deposit accounts represent hidden human trafficking financial risk.
		\item Customer E's former PEP status has fully expired after three years and has no impact on the risk score; the account should be treated as non-PEP.
	\end{enumerate}
	\textbf{Correct Answers: A and B}
	
	%--- Q62 ---
	\subsection{62. Institutional SAR Decision: Filing vs.\ Non-Filing Consequences \textnormal{[SELECT ONE]}}
	
	A compliance officer has completed an investigation into a corporate account and has reached the conclusion that there is reasonable suspicion of money laundering. The compliance committee is divided: three members believe a SAR should be filed; two members (including the revenue-generating relationship manager) argue that the evidence is ``not conclusive'' and that filing would ``damage the relationship.'' The BSA Officer decides to support the non-filing position, citing the lack of conclusive proof.
	
	Which of the following most accurately describes the legal and regulatory risk of the BSA Officer's decision?
	
	\begin{enumerate}[label=\Alph*.]
		\item The BSA Officer's decision is legally sound. SAR filing requires conclusive evidence of money laundering; reasonable suspicion is an insufficient legal standard.
		\item The BSA Officer's decision creates material regulatory risk. The SAR filing standard in the U.S. (and under FATF standards globally) is ``reasonable suspicion'' -- not certainty or conclusive proof. A compliance officer who concludes there is reasonable suspicion of ML has a legal obligation to file. Failure to file when the reasonable suspicion standard is met exposes the institution to regulatory sanctions (including civil money penalties), the BSA Officer to personal liability, and potentially the institution to criminal prosecution. Commercial relationship considerations are never a legally permissible basis for non-filing.
		\item The BSA Officer can resolve the situation by filing a CTR instead of a SAR, which is a less disruptive reporting mechanism.
		\item The three committee members who favored filing should be overruled because BSA Officers have unilateral authority to decide SAR filing questions.
		\item The institution should file the SAR but notify the customer simultaneously to maintain the commercial relationship.
	\end{enumerate}
	\textbf{Correct Answer: B}
	
	%--- Q63 ---
	\subsection{63. Proliferation Financing: Sectoral Sanctions Application \textnormal{[SELECT ONE]}}
	
	A Canadian bank is processing a trade finance transaction for a Russian energy company that is not individually listed on the OFAC SDN list or the Canadian SEMA list. However, the entity is subject to EU sectoral sanctions (OJ Regulation 833/2014) due to its involvement in deep-water Arctic oil exploration activities. The transaction involves a five-year term loan of USD 200 million, which would finance new drilling equipment. The bank's compliance officer notes that because Canada is not the EU, EU sectoral sanctions technically do not apply to Canadian entities.
	
	What is the compliance officer's most significant analytical error?
	
	\begin{enumerate}[label=\Alph*.]
		\item The compliance officer is correct; EU sectoral sanctions are strictly EU-law instruments and have no legal force outside EU member states.
		\item The compliance officer has failed to consider two critical dimensions: (1) because the transaction is denominated in USD, it may pass through U.S. correspondent banks that are subject to OFAC's Russia-related energy sector restrictions under Executive Orders 13662 and 14024, exposing the correspondent bank and potentially implicating the Canadian bank in a sanctions violation; and (2) even absent direct legal application, Canadian banks processing transactions connected to EU-sanctioned activities face significant reputational, correspondent banking, and de facto secondary sanctions risk that must be assessed and documented under a risk-based approach. The compliance officer should escalate for legal review and cannot simply rely on the technical non-application of EU law in Canada.
		\item The transaction is permissible because the entity is not on the SDN list, which is the only sanctions list that has extra-territorial effect.
		\item The five-year maturity of the loan is the primary concern, not the sanctions exposure; Canada's own sanctions framework restricts long-term financing for Russian entities.
		\item Canadian banks are fully insulated from U.S. secondary sanctions by the CUSMA (Canada-U.S.-Mexico Agreement) trade framework.
	\end{enumerate}
	\textbf{Correct Answer: B}
	
	%--- Q64 ---
	\subsection{64. Financial Investigation: Following the Money Through Asset Layers \textnormal{[SELECT TWO]}}
	
	A multi-agency financial crimes task force is building a money laundering case against a criminal organization. The investigative team has traced the following financial path of \$35 million in drug trafficking proceeds:
	
	\begin{enumerate}
		\item Cash deposited in sub-\$10,000 increments into 60 accounts at 12 different banks over 30 months.
		\item Funds consolidated into 8 shell company accounts and then used to purchase a portfolio of 22 residential properties.
		\item Rental income from the 22 properties deposited into a single corporate LLC account for 5 years.
		\item LLC assets liquidated and proceeds invested in a diversified publicly listed securities portfolio.
		\item Securities portfolio pledged as collateral for a \$40 million bank loan, which is used to fund a legitimate manufacturing business acquisition.
	\end{enumerate}
	
	Which two of the following statements are correct about the investigation's asset tracing challenges?
	
	\begin{enumerate}[label=\Alph*.]
		\item The commingling of criminal proceeds (drug cash) with legitimate rental income at Step (3) creates a significant asset tracing challenge, because investigators must distinguish the criminal portion of the LLC's funds from legitimately generated rental revenue to support a proportionate forfeiture order.
		\item By Step (5), the proceeds have been fully integrated. The bank loan, secured against a securities portfolio that originated from criminal proceeds, extends the financial crime into an otherwise legitimate business acquisition, potentially tainting the manufacturing business and the loan collateral with proceeds of crime -- both of which could be subject to forfeiture even though the manufacturing business itself is legitimate.
		\item Once proceeds are converted into real estate (Step 2), they are no longer traceable and the investigative trail ends; forfeiture cannot reach assets that have been converted into a different form.
		\item The five-year rental income period (Step 3) automatically legitimizes the funds under the U.S. statute of limitations for money laundering.
		\item The bank that made the Step (5) loan is automatically criminally liable for money laundering because it accepted criminal proceeds as collateral.
	\end{enumerate}
	\textbf{Correct Answers: A and B}
	
	%--- Q65 ---
	\subsection{65. The Risk-Based Approach: Calibrating EDD Intensity \textnormal{[SELECT ONE]}}
	
	Under FATF and CAMS standards, Enhanced Due Diligence (EDD) is not a fixed checklist but a flexible, risk-proportionate set of measures that must be scaled to the specific risk profile of the customer, product, and transaction. A compliance officer is designing EDD procedures for three distinct high-risk scenarios:
	
	\begin{itemize}
		\item Scenario X: A domestic PEP (provincial-level government official) opening a standard personal savings account with monthly deposits of \$5,000 from salary.
		\item Scenario Y: A foreign PEP (senior minister of a high-corruption-risk country) establishing a private banking relationship to manage a \$45 million portfolio funded from ``historical business income'' in the same sector as his ministry.
		\item Scenario Z: A foreign corporate client from a FATF Black List jurisdiction applying for a \$2 million trade finance facility for the import of dual-use electronic components.
	\end{itemize}
	
	Which statement best describes the correct differentiation of EDD intensity across these three scenarios?
	
	\begin{enumerate}[label=\Alph*.]
		\item All three scenarios require identical EDD procedures because all are classified as high risk.
		\item Scenario X requires proportionate EDD: verification of PEP status, source of salary funds, and board-level approval; the risk is lower given the domestic jurisdiction and low transaction value. Scenario Y requires the most intensive EDD: full source of wealth investigation, corroborated documentation of business income, senior management sign-off, and ongoing transaction-by-transaction review; the combination of foreign PEP status, high-corruption jurisdiction, large portfolio, and ministry-sector-linked funds creates the highest risk profile. Scenario Z requires EDD focused specifically on the dual-use goods' end-user certification, the trade finance documentation chain, the economic rationale for the transaction, and UNSC/OFAC sanctions screening for the Black List jurisdiction; this EDD is transaction-specific rather than relationship-wide. The RBA requires calibrating the depth of EDD to the specific risk drivers present in each scenario.
		\item Scenario Z requires no EDD because it involves a corporate client rather than a PEP.
		\item Scenarios X and Y require EDD; Scenario Z requires only Standard CDD because it is a trade finance product, not a wealth management relationship.
		\item EDD cannot be differentiated; regulatory consistency requires applying the same maximum-intensity EDD to all three scenarios to avoid any compliance gap.
	\end{enumerate}
	\textbf{Correct Answer: B}
	
	%--- Q66 ---
	\subsection{66. Correspondent Banking: Termination Protocol and De-Risking Documentation \textnormal{[SELECT ONE]}}
	
	A bank has decided, following a documented individual risk assessment, to terminate its correspondent banking relationship with a respondent bank in a high-risk jurisdiction. The relationship has generated \$1.2 million in annual fees and the respondent bank has maintained satisfactory regulatory standing. The termination decision is based on the risk profile of the respondent bank's underlying customer base (heavy concentration of unlicensed MSBs and high-risk NBFI activity) and the correspondent's inability to provide adequate documentation on downstream clients after multiple formal RFIs. The respondent bank's management calls the correspondent's CEO, arguing that the termination is discriminatory and that the correspondent bank is engaging in illegal de-risking.
	
	Which of the following most accurately describes the correspondent bank's legal and regulatory position?
	
	\begin{enumerate}[label=\Alph*.]
		\item The respondent bank's argument is correct: a bank with satisfactory regulatory standing cannot legally be terminated for AML risk management reasons, because this constitutes discriminatory de-risking.
		\item The correspondent bank's decision is legally defensible and consistent with FATF and Wolfsberg guidance. Banks are not required to maintain relationships that present unacceptable AML risk. The termination decision was based on a documented individual risk assessment (not a blanket categorical exclusion), specifically identifying the downstream customer base composition and the respondent's failure to respond adequately to multiple RFIs -- both of which are legitimate risk-based grounds for termination. This is not indiscriminate de-risking; it is a justified risk-based exit from a relationship that cannot be adequately monitored.
		\item The bank must first obtain a regulatory no-objection letter before terminating a long-standing correspondent relationship.
		\item The termination must be reversed because the respondent bank's satisfactory regulatory standing demonstrates adequate AML controls, which overrides the correspondent's own risk assessment.
		\item The correspondent bank must provide the respondent bank with a detailed explanation of every specific SAR filed on the account as justification for the termination.
	\end{enumerate}
	\textbf{Correct Answer: B}
	
	%--- Q67 ---
	\subsection{67. Sanctions: OFAC Voluntary Self-Disclosure and Mitigation \textnormal{[SELECT TWO]}}
	
	A U.S. bank's compliance team discovers during an internal audit that its sanctions screening system failed to flag 11 wire transfers processed over 18 months that involved an entity on the OFAC SDN list. The failure was due to a known software deficiency that had been flagged in a vendor audit but not remediated. The total value of the 11 transactions was \$430,000. The bank's legal team is evaluating whether to submit a Voluntary Self-Disclosure (VSD) to OFAC.
	
	Which two of the following statements correctly describe OFAC's VSD process and its consequences?
	
	\begin{enumerate}[label=\Alph*.]
		\item Filing a Voluntary Self-Disclosure with OFAC is a mitigating factor that can significantly reduce the civil monetary penalty imposed for the sanctions violations. OFAC's Enforcement Guidelines provide that a VSD can result in a 50\% reduction of the base penalty.
		\item A VSD requires the bank to provide a full factual account of the violations, including the root cause analysis (in this case, the unresolved software deficiency), the transactions affected, the total value, and a remediation plan. A thorough and accurate VSD that demonstrates genuine compliance culture and remediation commitment is treated more favorably by OFAC.
		\item Filing a VSD automatically triggers a criminal referral to the Department of Justice; therefore banks should never self-disclose.
		\item A VSD is only available for violations involving amounts under \$1,000; the \$430,000 aggregate value disqualifies the bank from submitting a VSD.
		\item OFAC treats the pre-existing knowledge of the software deficiency as an aggravating factor regardless of whether the bank files a VSD.
	\end{enumerate}
	\textbf{Correct Answers: A and B}
	
	%--- Q68 ---
	\subsection{68. Enforcement Action: Identifying Root Cause of a BSA Program Failure \textnormal{[SELECT ONE]}}
	
	A federal banking regulator has issued a Consent Order against a mid-sized bank after a multi-year examination. The Order cites the following findings: (1) transaction monitoring system thresholds were calibrated only once at implementation five years ago and never revised; (2) the BSA Officer position was vacant for 11 months before being filled with an unqualified candidate; (3) independent testing was conducted by the same team that designed the monitoring scenarios; (4) KYC files for 34\% of high-risk accounts were missing required EDD documentation; and (5) SAR narratives for 60\% of filed SARs lacked sufficient detail to enable law enforcement to use the reports. The regulator requires a comprehensive remediation plan.
	
	What is the most accurate characterization of the root cause underlying all five findings?
	
	\begin{enumerate}[label=\Alph*.]
		\item The root cause is a technology deficiency: the transaction monitoring system is outdated and requires replacement.
		\item The five findings collectively indicate a systemic governance and organizational failure. The BSA program has suffered from a lack of senior management commitment and board oversight (allowing a key position to remain vacant for 11 months), inadequate resource allocation (insufficient staffing for EDD and SAR quality review), institutional independence failures (audit team self-reviewing its own scenarios), and no program maintenance culture (monitoring never recalibrated in five years). These are enterprise-wide governance deficiencies that cannot be resolved by technology upgrades alone; they require structural changes to accountability, resources, and board-level engagement with the AML program.
		\item Each finding is an isolated operational incident; no single root cause unifies the five deficiencies.
		\item The root cause is the regulatory examination process itself, which applied standards that were not in place when the bank originally designed its AML program.
		\item The root cause is insufficient SAR volume; the bank should increase the number of SARs filed to address all five findings.
	\end{enumerate}
	\textbf{Correct Answer: B}
	
	%--- Q69 ---
	\subsection{69. Public-Private Information Sharing: FinCEN Section 314(b) \textnormal{[SELECT ONE]}}
	
	Two banks, NorthBank and SouthBank, each have accounts held by different individuals who appear, based on their respective internal investigations, to be part of the same money laundering network. Neither bank's internal investigation alone has been sufficient to establish the complete picture of the network. NorthBank's BSA Officer has heard that SouthBank may have related account information and is considering whether to contact SouthBank directly to share suspicious activity information.
	
	Under FinCEN's Section 314(b) program, which of the following accurately describes the permissible scope of information sharing?
	
	\begin{enumerate}[label=\Alph*.]
		\item NorthBank and SouthBank cannot share SAR-related information under any circumstances because all SAR information is subject to absolute confidentiality and may only be shared with the FIU.
		\item Under Section 314(b), financial institutions that have filed the required FinCEN opt-in notification may voluntarily share information with other opted-in institutions about persons they suspect of money laundering or terrorist financing activity. The sharing institutions receive a safe harbor from liability for the sharing. NorthBank and SouthBank may exchange information about the suspected network to build a more complete picture, provided both institutions have filed their 314(b) notifications and the sharing is for the purpose of identifying and reporting ML/TF activity.
		\item Section 314(b) requires both institutions to first obtain a court order before exchanging customer information, even for AML purposes.
		\item Section 314(b) sharing is only permissible if the amount of suspected money laundering exceeds \$1 million; smaller suspected ML activity must be reported individually by each institution without inter-institution sharing.
		\item Section 314(b) sharing automatically constitutes tipping-off because the shared information relates to a suspected subject of an AML investigation.
	\end{enumerate}
	\textbf{Correct Answer: B}
	
	%--- Q70 ---
	\subsection{70. Comprehensive Final Scenario: The Cross-Border Corruption and Integration Scheme \textnormal{[SELECT TWO]}}
	
	\textbf{Background:} Over seven years, a network involving a foreign government minister, a private equity fund manager, and a real estate developer has laundered approximately \$85 million in bribery proceeds originating from state infrastructure contracts in Country P. The mechanism: (1) bribe payments in USD cash are collected in Country P and structured into 200+ accounts at 18 different banks across four countries; (2) funds are consolidated into a Cayman Islands private equity fund (Fund Alpha) managed by an investment manager in London; (3) Fund Alpha invests in commercial real estate in three European capitals; (4) the properties generate legitimate rental income and are sold after five years at a combined gain of \$25 million; (5) the sale proceeds and rental income are distributed to the fund's limited partners, which are shell companies ultimately controlled by the minister and his associates.
	
	A newly appointed compliance officer at the London investment manager conducts a routine file review and identifies three red flags: (a) the fund's limited partners are all offshore shell companies with no verifiable business purpose; (b) the fund's capital commitments were funded by wire transfers from accounts in jurisdictions known for lax AML; (c) the investment manager has never conducted source of funds or beneficial ownership verification on the LP entities.
	
	\textbf{Question:} Which two of the following represent both the most critical immediate action and the most accurate characterization of the investment manager's compliance failures?
	
	\begin{enumerate}[label=\Alph*.]
		\item The most critical immediate action is to suspend new capital calls from the LP entities, conduct emergency EDD on all LP entities to establish UBO and source of funds, escalate the findings to senior management and the compliance committee, and file an STR with the UK National Crime Agency (NCA) if the findings support reasonable suspicion that the LP investments represent criminal proceeds.
		\item The investment manager has committed AML compliance failures by (i) onboarding LP entities without conducting UBO identification and source of funds verification, violating basic CDD obligations applicable to investment managers under UK AML regulations; and (ii) failing to conduct ongoing monitoring sufficient to identify that the LP capital originated from accounts in high-risk jurisdictions without any documented economic rationale.
		\item The investment manager has no AML obligation for LP investors because private equity funds are not regulated financial institutions under the EU Anti-Money Laundering Directives.
		\item The correct action is to return all LP capital to the originating accounts immediately and close Fund Alpha, because the investment manager cannot be held responsible for the prior failures of the 18 banks that processed the original structured deposits.
		\item The three red flags identified by the compliance officer are not sufficient to trigger STR filing; a SAR can only be filed after a conviction has been secured against the minister.
	\end{enumerate}
	\textbf{Correct Answers: A and B}
	
\end{document}